%% 本文件原为 工程文件/scripts/assemble.py 生成的 \input 骨架。
%% 第二轮校订（2026-11-21）为直接修改译文而把 units/ch02_u01--u04 内联展开；
%% 内容与原装配结果等价。请勿再运行 assemble.py，否则会覆盖校订。
\bdchap{ch:quant}{Hitchin 哈密顿量的量子化}
% ch02_u01: 原书 p.15--p.23
\SourcePage{15}

\bdsec{2.1. $\operatorname{Bun}_G$ 的几何}

我们沿用 0.1 的记号；特别地，$G$ 是半单群，$X$ 是亏格 $g>1$ 的光滑射影曲线。

\bdp{2.1.1} 众所周知，$\operatorname{Bun}_G$ 是纯维数 $(g-1)\dim G$ 的光滑代数栈。$\operatorname{Bun}_G$ 的连通分支集合可（经「第一陈类」映射）典范地等同于 $H^{2}(X,\pi_{1}^{et}(G))=\pi_{1}(G)$。这里 $\pi_{1}^{et}(G)$ 是 Grothendieck 意义下的基本群；$\pi_{1}(G)$ 是 $G$ 的余权群模去余根子群所得的商。二者相差一个 Tate 扭：$\pi_{1}^{et}(G)=\pi_{1}(G)(1)$。

对 $\mathcal{F}\in\operatorname{Bun}_G$，切层 $\Theta=\Theta_{\operatorname{Bun}_G}$ 在 $\mathcal{F}$ 处的纤维是 $H^{1}(X,\mathfrak{g}_{\mathcal{F}})$。说明一下：对 $G$-模 $W$，记 $W_{\mathcal{F}}$ 为 $W$ 的 $\mathcal{F}$-扭，它是 $X$ 上的向量丛。这里经伴随作用把 $\mathfrak{g}$ 看作 $G$-模。

按定义，典范线丛 $\omega=\omega_{\operatorname{Bun}_G}$ 是 $\operatorname{Bun}_G$ 的余切复形的行列式（见 [LMB93]）。该复形在 $\mathcal{F}\in\operatorname{Bun}_G$ 上的纤维与 $R\Gamma(X,\mathfrak{g}_{\mathcal{F}})[1]$ 对偶（见 [LMB93]），故 $\omega$ 在 $\mathcal{F}$ 上的纤维是 $\det R\Gamma(X,\mathfrak{g}_{\mathcal{F}})$。\footnote{作者们不该忘记核验 [LMB93] 中确实包含此处所断言的内容！！}

\bdp{2.1.2} \emph{命题.} $\operatorname{Bun}_G$ 在 1.1.1 的意义下是极良的。

证明将在 2.10.5 给出。事实上，我们将用到 $\operatorname{Bun}_G$ 是良的这一事实。由 1.1，有泊松代数层 $P=P_{\operatorname{Bun}_G}$ 与扭微分算子层 $D^{\lambda}=D_{\operatorname{Bun}_G,\omega^{\lambda}}$。众所周知，当 $\lambda\neq 1/2$ 时，$D^{\lambda}$ 的整体截面只有局部常值函数。在第 2 节与第 3 节中我们处理 $D':=D^{1/2}$；其截面就简称为扭微分算子。

\bdsec{2.2. Hitchin 构造 I}

\SourcePage{16}

\bdp{2.2.1} 置 $\mathcal{C}=\mathcal{C}_{\mathfrak{g}}:=\operatorname{Spec}(\operatorname{Sym}\mathfrak{g})^{G}$；这是 $\mathfrak{g}^{*}$ 关于余伴随作用的仿射概形商。$\mathcal{C}$ 带有乘法群 $\mathbb{G}_m$ 的典范作用，它来自 $\mathfrak{g}^{*}$ 上的位似变换。非典范地取一组次数为 $d_i$、$i\in I$ 的齐次生成元 $p_i\in(\operatorname{Sym}\mathfrak{g})^{G}$，就把 $\mathcal{C}$ 与坐标空间 $\mathbb{C}^{I}$ 等同；元素 $\lambda\in\mathbb{G}_m$ 通过对角矩阵 $(\lambda^{d_i})$ 作用。

\bdp{2.2.2} 记 $\mathcal{C}_{\omega_X}$ 为 $\mathcal{C}$ 关于上述 $\mathbb{G}_m$-作用的 $\omega_X$-扭（这里把典范丛 $\omega_X$ 看作 $X$ 上的 $\mathbb{G}_m$-torsor，即主齐性空间）。它是 $X$ 上的丛；上述 $p_i$ 把 $\mathcal{C}_{\omega_X}$ 与 $\prod_I\omega_X^{\otimes d_i}$ 等同。置
\[
\operatorname{Hitch}(X)=\operatorname{Hitch}_{\mathfrak{g}}(X):=\Gamma(X,\mathcal{C}_{\omega_X}).
\]
换言之，$\operatorname{Hitch}(X)=\operatorname{Mor}((\operatorname{Sym}^{\cdot}\mathfrak{g})^{G},\Gamma(X,\omega_X^{\otimes\cdot}))$（分次代数的态射）。把 $\operatorname{Hitch}(X)$ 看作带 $\mathbb{G}_m$-作用的代数簇；它非典范地同构于向量空间 $\prod_I\Gamma(X,\omega_X^{\otimes d_i})$。存在唯一的点 $0\in\operatorname{Hitch}(X)$，它在 $\mathbb{G}_m$ 的作用下不动。记 $\mathfrak{z}^{cl}(X)=\mathfrak{z}_{\mathfrak{g}}^{cl}(X)$ 为 $\operatorname{Hitch}(X)$ 上的函数环；这是分次交换代数。更确切地说，$\mathfrak{z}^{cl}(X)$ 的分次对应于 $\mathfrak{z}^{cl}(X)$ 上一个 $\mathbb{G}_m$-作用，它与 $\mathcal{C}$ 上 $\mathbb{G}_m$-作用诱导的作用相反；因此 $\mathfrak{z}^{cl}(X)$ 的分次为正。

\bdp{2.2.3} 由 Serre 对偶与 2.1.1，$\mathcal{F}\in\operatorname{Bun}_G$ 处 $\operatorname{Bun}_G$ 的余切空间 $T_{\mathcal{F}}^{*}\operatorname{Bun}_G$ 等于 $\Gamma(X,\mathfrak{g}_{\mathcal{F}}^{*}\otimes\omega_X)$。$G$-不变投影 $\mathfrak{g}^{*}\to\mathcal{C}$ 给出态射 $\mathfrak{g}_{\mathcal{F}}^{*}\otimes\omega_X\to\mathcal{C}_{\omega_X}$，以及映射 $p_{\mathcal{F}}:T_{\mathcal{F}}^{*}\operatorname{Bun}_G\to\operatorname{Hitch}(X)$。当 $\mathcal{F}$ 变动时得到态射
\[
p:T^{*}\operatorname{Bun}_G\to\operatorname{Hitch}(X)
\]
或等价地，分次交换代数的态射
\[
h_X^{cl}:\mathfrak{z}^{cl}(X)\to\Gamma(T^{*}\operatorname{Bun}_G,\mathcal{O})=\Gamma(\operatorname{Bun}_G,P).
\]
$p$ 称为 Hitchin 纤维化。

\SourcePage{17}

记 $\operatorname{Bun}_G^{\gamma}$ 为 $\operatorname{Bun}_G$ 中对应于 $\gamma\in\pi_1(G)$ 的连通分支（见 2.1.1），记 $p^{\gamma}$ 为 $p$ 到 $T^{*}\operatorname{Bun}_G^{\gamma}$ 的限制。

\bdp{2.2.4} \emph{定理.}（[Hit87]，[Fal93]，[Gi97]）
\begin{enumerate}[label=(\roman*),leftmargin=2.6em]
\item $h_X^{cl}$ 的像由泊松交换的函数组成。
\item $\dim\operatorname{Hitch}(X)=\dim\operatorname{Bun}_G=(g-1)\cdot\dim\mathfrak{g}$。
\item $p$ 是平坦的，其纤维的维数纯为 $\dim\operatorname{Bun}_G$。对每个 $\gamma\in\pi_1(X)$，$p^{\gamma}$ 是满射。
% [Tutor Note: 原书 p.17 定理 2.2.4 的 (iii) 与 (v) 均印作 γ ∈ π₁(X)，按上下文应为 π₁(G)；底本如此，保留原样。]
\item 存在非空开集 $U\subset\operatorname{Hitch}(X)$，使得对任意 $\gamma\in\pi_1(G)$，态射 $(p^{\gamma})^{-1}(U)\to U$ 固有且光滑，并且其纤维连通。事实上，$p^{\gamma}$ 在 $u\in U$ 上的纤维同构于某个阿贝尔簇 $A_u$ 与中心 $Z\subset G$ 的分类栈之积。
\item 对每个 $\gamma\in\pi_1(X)$，态射 $\mathfrak{z}^{cl}(X)\to\Gamma(\operatorname{Bun}_G^{\gamma},P)$ 是同构。\hfill$\Box$
\end{enumerate}

\emph{注.}
\begin{enumerate}[label=(\roman*),leftmargin=2.6em]
\item 毋庸赘言，定理 2.2.4 的主要贡献属于 Hitchin [Hit87]。
\item 定理 2.2.4 蕴含 $p$ 是 Lagrange 纤维化；或者，若你愿意这么说，来自 $h_X^{cl}(\mathfrak{z}^{cl}(X))$ 的哈密顿量在 $T^{*}\operatorname{Bun}_G$ 上定义完全可积系统。在栈的语境中我们不怕用这些词，因为 Lagrange 纤维化的概念是双有理的。又因 $\operatorname{Bun}_G$ 在 1.1.1 的意义下是极良的，$T^{*}\operatorname{Bun}_G$ 有开稠密的 Deligne--Mumford 子栈 $T^{*}\operatorname{Bun}_G^{0}$，它在显然的意义下是辛的（这里 $\operatorname{Bun}_G^{0}$ 是自同构群有限的 $G$-丛构成的栈）。
\item Hitchin 在 [Hit87] 中给出陈述 (i) 的复分析证明。我们将在 2.4.3 给出 (i) 的代数证明。
\SourcePage{18}
\item Hitchin 对 (ii) 的证明很简单：由 2.2.2，$\dim\operatorname{Hitch}(X)=\sum_i\dim\Gamma(X,\omega_X^{\otimes d_i})$；因 $g>1$，$\dim\Gamma(X,\omega_X^{\otimes d_i})=(g-1)(2d_i-1)$；最后 $(g-1)\sum_i(2d_i-1)=(g-1)\dim\mathfrak{g}=\dim\operatorname{Bun}_G$。
\item 对经典群 $G$，陈述 (iv) 已由 Hitchin [Hit87] 证明。一般情形由 Faltings 证明（[Fal93] 中的定理 III.2）。
\item 陈述 (v) 由 (iii) 与 (iv) 推出。
\item (iii) 的证明的一些评述将在 2.10 给出。
\end{enumerate}

\bdp{2.2.5} 我们的目标是解决如下量子化问题：构造滤交换代数 $\mathfrak{z}(X)$，并赋予它同构 $\sigma_{\mathfrak{z}(X)}:\operatorname{gr}\mathfrak{z}(X)\xrightarrow{\sim}\mathfrak{z}^{cl}(X)$，以及滤代数的态射 $h_X:\mathfrak{z}(X)\to\Gamma(\operatorname{Bun}_G,D')$；后者与象征映射相容，即 $\sigma_{\operatorname{Bun}_G}\circ\operatorname{gr}h_X=h_X^{cl}\circ\sigma_{\mathfrak{z}(X)}$（$\sigma_{\operatorname{Bun}_G}$ 的定义见 1.1.4 与 1.1.6）。

注意，2.2.4(v) 于是蕴含：对任意 $\gamma\in\pi_1(X)$，映射 $h_X^{\gamma}:\mathfrak{z}(X)\to\Gamma(\operatorname{Bun}_G^{\gamma},D')$ 是同构。
% [Tutor Note: 原书 p.18 此处亦印作 γ ∈ π₁(X)，按上下文应为 π₁(G)；底本如此，保留原样。]
因此若 $G$ 单连通，则这样的构造唯一；它归结为两条断言：$\Gamma(\operatorname{Bun}_G,D')$ 是交换代数；$T^{*}\operatorname{Bun}_G$ 上的任一整体函数都是某个整体扭微分算子的象征。

我们不知道如何用整体考虑直接解决此问题。我们将从 Hitchin 图景的局部版本出发，遵循 1.2 的量子化方案。上述量子化问题的同一个解将给出两个构造。第一个（见 2.5.5）较易表述；第二个（见 2.7.4）的优点是完全典范。为证明第一个构造确实给出解，我们使用第二个。正是第二个构造将给出 $\operatorname{Spec}\mathfrak{z}(X)$ 与 $X$ 上 $(^{L}G)_{ad}$-局部系统栈的某个子空间的等同（见 3.3.2）。

\bdsec{2.3. $\operatorname{Bun}_G$ 的几何 II}

让我们回忆 $\operatorname{Bun}_G$ 如何纳入 1.2.6 的框架。

\SourcePage{19}

\bdp{2.3.1} 固定一点 $x\in X$。记 $\mathcal{O}_x$ 为 $x$ 的完备局部环，$K_x$ 为其分式域。设 $\mathfrak{m}_x\subset\mathcal{O}_X$ 为极大理想。置 $\mathcal{O}_x^{(n)}:=\mathcal{O}_X/\mathfrak{m}_x^{n}$（于是 $\mathcal{O}_x=\varprojlim\mathcal{O}_x^{(n)}$）。群 $G(\mathcal{O}_x^{(n)})$ 是某个仿射代数群的 $\mathbb{C}$-点群；我们仍把它记作 $G(\mathcal{O}_x^{(n)})$（滥用记号）。$G(\mathcal{O}_x^{(n)})$ 是 $G(\mathcal{O}_x^{(n+1)})$ 的商。于是 $G(\mathcal{O}_x)=\varprojlim G(\mathcal{O}_x^{(n)})$ 是仿射群概形。

记 $\operatorname{Bun}_{G,nx}$ 为 $X$ 上在 $\operatorname{Spec}\mathcal{O}_x^{(n)}$ 上平凡化的 $G$-丛构成的栈（注意除子 $nx$ 与子概形 $\operatorname{Spec}\mathcal{O}_x^{(n)}\subset X$ 相同）。这是 $\operatorname{Bun}_G$ 上的 $G(\mathcal{O}_x^{(n)})$-torsor。把 $\operatorname{Bun}_{G,nx}$ 的点记作 $(\mathcal{F},\alpha^{(n)})$。有显然的仿射投影 $\operatorname{Bun}_{G,(n+1)x}\to\operatorname{Bun}_{G,nx}$。置 $\operatorname{Bun}_{G,\underline{x}}:=\varprojlim\operatorname{Bun}_{G,nx}$；这是 $\operatorname{Bun}_G$ 上的 $G(\mathcal{O}_x)$-torsor。

\bdp{2.3.2} \emph{命题.} $\operatorname{Bun}_{G,\underline{x}}$ 是概形。$\operatorname{Bun}_{G,\underline{x}}$ 上的 $G(\mathcal{O}_x)$-作用满足 1.2.6 的条件 (16)。\hfill$\Box$

\bdp{2.3.3} 众所周知，$\operatorname{Bun}_{G,\underline{x}}$ 上的 $G(\mathcal{O}_x)$-作用可典范地延拓为群归纳概形 $G(K_x)$ 的作用（归纳概形的定义见 7.11.1，归纳概形 $G(K_x)$ 的定义见 7.11.2 (iv)）。因 $\operatorname{Lie}G(K_x)=\mathfrak{g}\otimes K_x$，特别地有 Harish-Chandra 对 $(\mathfrak{g}\otimes K_x,G(\mathcal{O}_x))$ 在 $\operatorname{Bun}_{G,\underline{x}}$ 上的作用。

让我们回忆 $G(K_x)$-作用的定义。由 7.11.2 (iv)，须对任意 $\mathbb{C}$-代数 $R$ 定义 $G(R\widehat{\otimes}K_x)$ 在 $\operatorname{Bun}_{G,\underline{x}}(R)$ 上的作用。为此使用下述定理，它本质上属于 A. Beauville 与 Y. Laszlo。置 $X':=X\setminus\{x\}$。

\bdp{2.3.4} \emph{定理.} $X\otimes R$ 上的 $G$-丛 $\mathcal{F}$ 等同于三元组 $(\mathcal{F}_1,\mathcal{F}_2,\varphi)$：其中 $\mathcal{F}_1$ 是 $X'\otimes R$ 上的 $G$-丛，$\mathcal{F}_2$ 是 $\operatorname{Spec}(R\widehat{\otimes}\mathcal{O}_x)$ 上的 $G$-丛，$\varphi$ 是 $\mathcal{F}_1$ 与 $\mathcal{F}_2$ 到 $\operatorname{Spec}(R\widehat{\otimes}K_x)$ 的拉回之间的同构。

更确切地说，从 $X\otimes R$ 上 $G$-丛 $\mathcal{F}$ 的范畴（=群胚）到上述三元组 $(\mathcal{F}_1,\mathcal{F}_2,\varphi)$ 的范畴的函子，由 $\mathcal{F}_1:=\mathcal{F}|_{X'\otimes R}$、$\mathcal{F}_2:=\mathcal{F}$ 到 $\operatorname{Spec}(R\widehat{\otimes}\mathcal{O}_x)$ 的拉回、$\varphi:=\mathrm{id}$ 定义；它是范畴等价。

\SourcePage{20}

由该定理，$\operatorname{Bun}_{G,\underline{x}}$ 的 $R$-点等同于 $X'\otimes R$ 上的一个 $G$-丛，连同一个其到 $\operatorname{Spec}(R\widehat{\otimes}K_x)$ 的拉回的平凡化。于是 $G(R\widehat{\otimes}K_x)$ 通过改变平凡化作用在 $\operatorname{Bun}_{G,\underline{x}}(R)$ 上。这样就得到 $G(K_x)$ 在 $\operatorname{Bun}_{G,\underline{x}}$ 上的作用。

定理 2.3.4 的证明基于下述定理，它是 [BLa95] 主要结果的一个特例。

\bdp{2.3.5} \emph{定理.}（Beauville--Laszlo）平坦拟凝聚 $\mathcal{O}_{X\otimes R}$-模 $\mathcal{M}$ 的范畴等价于三元组 $(\mathcal{M}_1,\mathcal{M}_2,\varphi)$ 的范畴。其中 $\mathcal{M}_1$ 是 $X'\otimes R$ 上的平坦拟凝聚 $\mathcal{O}$-模，$\mathcal{M}_2$ 是 $\operatorname{Spec}(R\widehat{\otimes}\mathcal{O}_x)$ 上的平坦拟凝聚 $\mathcal{O}$-模，$\varphi$ 是 $\mathcal{M}_1$ 与 $\mathcal{M}_2$ 到 $\operatorname{Spec}(R\widehat{\otimes}K_x)$ 的拉回之间的同构（从第一个范畴到第二个范畴的函子按定理 2.3.4 定义）。$\mathcal{M}$ 局部自由且秩有限，当且仅当相应的 $\mathcal{M}_1$ 与 $\mathcal{M}_2$ 具有此性质。

\emph{注.} 若 $R$ 诺特且层凝聚，则有 M. Artin 的一个远为一般的「粘合定理」（[Ar] 中的定理 2.6）。但因 $G(K_x)$ 的子概形通常是无限型的，我们使用 Beauville--Laszlo 定理，它无需诺特假设。

为从 2.3.5 推出定理 2.3.4，只需把 $G$-丛解释为张量函子 $\{G\text{-模}\}\to\{\text{向量丛}\}$。或者，可把 $X\otimes R$ 上的 $G$-丛解释为主 $G$-丛，即平坦仿射态射 $\pi:\mathcal{F}\to X\otimes R$，连同 $G$ 在 $\mathcal{F}$ 上满足某些性质的作用；于是可用层 $\mathcal{M}:=\pi_{*}\mathcal{O}_{\mathcal{F}}$ 改写这些数据，并应用定理 2.3.5。
% [Tutor Note: 原书 p.20 此处印作 "principle G-bundle"，应为 "principal G-bundle"；底本如此，译文按正字处理。]

\bdp{2.3.6} \emph{注.} 这里给出 $\mathfrak{g}\otimes K_x$ 在 $\operatorname{Bun}_{G,\underline{x}}$ 上作用的直接描述，该作用由 $G(K_x)$ 的作用诱导（将来不会用到 ???）。取 $(\mathcal{F},\bar\alpha)\in\operatorname{Bun}_{G,\underline{x}}$，$\bar\alpha=\varprojlim\alpha^{(n)}$。$\operatorname{Bun}_{G,nx}$ 在 $(\mathcal{F},\alpha^{(n)})$ 处的切空间是 $H^{1}(X,\mathfrak{g}_{\mathcal{F}}(-nx))$，故 $\Theta_{\operatorname{Bun}_{G,\underline{x}}}$ 在 $(\mathcal{F},\bar\alpha)$ 处的纤维等于 $\varprojlim H^{1}(X,\mathfrak{g}_{\mathcal{F}}(-nx))=H^{1}_{c}(X\setminus\{x\},\mathfrak{g}_{\mathcal{F}})$。有通常的满射 $\mathfrak{g}_{\mathcal{F}}\otimes_{\mathcal{O}_X}K_x\to H^{1}_{c}(X\setminus\{x\},\mathfrak{g}_{\mathcal{F}})$。用 $\bar\alpha$ 把 $\mathfrak{g}_{\mathcal{F}}\otimes_{\mathcal{O}_X}K_x$ 与 $\mathfrak{g}\otimes K_x$ 等同。
% [Tutor Note: 原书 p.20 此处交叉引用印作「???」，是 LaTeX 未解析的引用；底本如此，保留原样。]

\SourcePage{21}

当 $(\mathcal{F},\bar\alpha)$ 变动时得到映射 $\mathfrak{g}\otimes K_x\to\Theta_{\operatorname{Bun}_{G,\underline{x}}}$。我们的 $\mathfrak{g}\otimes K_x$-作用是该映射的相反数（???）。
% [Tutor Note: 原书 p.21 此处交叉引用印作「???」；底本如此，保留原样。]

\bdp{2.3.7} \emph{注.} 设 $D\subset X\otimes R$ 是 $\operatorname{Spec}R$ 上的有限闭子概形，它可局部由一个方程定义（即 $D$ 是有效相对 Cartier 除子）。记 $\tilde{D}$ 为 $D$ 的形式邻域，$A$ 为 $\tilde{D}$ 的坐标环（于是 $\tilde{D}$ 是仿射形式概形，$\operatorname{Spec}A$ 是真概形）。则把 $X'\otimes R$ 换成 $(X\otimes R)\setminus D$、把 $R\widehat{\otimes}\mathcal{O}_x$ 换成 $A$、把 $\operatorname{Spec}(R\widehat{\otimes}K_x)$ 换成 $(\operatorname{Spec}A)\setminus D$ 后，定理 2.3.4 与 2.3.5 仍然成立。

若 $D$ 的正规丛平凡，这由 [BLa95] 的主要定理推出：事实上，此时可构造仿射邻域 $U\supset D$，使得在 $U$ 内子概形 $D$ 由整体方程 $f=0$ 定义，$f\in H^{0}(U,\mathcal{O}_U)$（这是 [BLa95] 中考虑的情形）。\footnote{为构造 $U$ 与 $f$，注意当 $n$ 充分大时存在 $\varphi_n\in H^{0}(X\otimes R,\mathcal{O}_{X\otimes R}(nD))$，使得 $\mathcal{O}_{X\otimes R}(nD)/\mathcal{O}_{X\otimes R}((n-1)D)$ 由 $\varphi_n$ 生成；然后置 $U:=(X\otimes R)\setminus\{\varphi_n\varphi_{n+1}\text{ 的零点集}\}$，$f:=\varphi_n/\varphi_{n+1}$（此构造在 $D\to\operatorname{Spec}R$ 为满射时有效，这是一个无害的假设）。} 对本工作的目的而言，$D$ 的正规丛平凡的情形已足够。为处理一般情形，需要 [BLa95] 主要定理的整体化版本（见 2.12）。此外还需把态射 $\tilde{D}\to X\otimes R$ 延拓为态射 $\operatorname{Spec}A\to X\otimes R$（显然这样的延拓唯一，但其存在性须证明）；见 2.12。

\bdsec{2.4. Hitchin 构造 II}

\bdp{2.4.1} 置 $\omega_{\mathcal{O}_x}:=\varprojlim\omega_{\mathcal{O}_x^{(n)}}$，其中 $\omega_{\mathcal{O}_x^{(n)}}$ 是 $\mathcal{O}_x^{(n)}=\mathcal{O}_x/\mathfrak{m}_x^{n}$ 的微分模。记 $\operatorname{Hitch}_x^{(n)}$ 为 $\mathcal{C}_{\omega_X}$ 在 $\operatorname{Spec}\mathcal{O}_x^{(n)}$ 上的截面概形。这是带 $\mathbb{G}_m$-作用的仿射概形，非典范地同构于向量空间 $M/\mathfrak{m}_x^{n}M$，其中 $M:=\prod\omega_{\mathcal{O}_x}^{\otimes d_i}$。置
\[
\operatorname{Hitch}_x=\operatorname{Hitch}_{\mathfrak{g}}(\mathcal{O}_x):=\varprojlim\operatorname{Hitch}_x^{(n)}.
\]
这是带 $\mathbb{G}_m$-作用的仿射概形，非典范地同构于 $M=\prod\omega_{\mathcal{O}_x}^{\otimes d_i}$。于是 $\operatorname{Hitch}_x$ 是 $\mathcal{C}_{\omega_X}$ 在 $\operatorname{Spec}\mathcal{O}_x$ 上的截面概形。

\SourcePage{22}

记 $\mathfrak{z}_x^{cl}=\mathfrak{z}_{\mathfrak{g}}^{cl}(\mathcal{O}_x)$ 为 1.2.2 中的分次泊松代数 $P_{(\mathfrak{g}\otimes K_x,G(\mathcal{O}_x))}=\operatorname{Sym}(\mathfrak{g}\otimes K_x/\mathcal{O}_x)^{G(\mathcal{O}_x)}$。我们将构造典范的 $\mathbb{G}_m$-等变同构 $\operatorname{Spec}\mathfrak{z}_x^{cl}\xrightarrow{\sim}\operatorname{Hitch}_x$（$\mathfrak{z}_x^{cl}$ 上的 $\mathbb{G}_m$-作用与分次所诱导的作用相反；参看 2.2.2 末尾）。

留数配对把 $(K_x/\mathcal{O}_x)^{*}$ 与 $\omega_{\mathcal{O}_x}$ 等同，故 $\operatorname{Spec}\operatorname{Sym}(\mathfrak{g}\otimes K_x/\mathcal{O}_x)=\mathfrak{g}^{*}\otimes\omega_{\mathcal{O}_x}$。投影 $\mathfrak{g}^{*}\to\mathcal{C}$ 给出仿射概形的态射 $\mathfrak{g}^{*}\otimes\omega_{\mathcal{O}_x}\to\operatorname{Hitch}_x$。它 $G(\mathcal{O}_x)$-不变，故诱导态射 $\operatorname{Spec}\mathfrak{z}_x^{cl}\to\operatorname{Hitch}_x$。为说明这是同构，须证明：$\mathfrak{g}^{*}\otimes\omega_{\mathcal{O}_x}$ 上每个 $G(\mathcal{O}_x)$-不变正则函数都来自 $\operatorname{Hitch}_x$ 上唯一的一个正则函数。显然可把 $\mathfrak{g}^{*}\otimes\omega_{\mathcal{O}_x}$ 换成 $\mathfrak{g}^{*}\otimes\mathcal{O}_x=\operatorname{Paths}(\mathfrak{g}^{*})$，把 $\operatorname{Hitch}_x$ 换成 $\operatorname{Paths}(\mathcal{C})$（对概形 $Y$，记 $\operatorname{Paths}(Y)$ 为态射 $\operatorname{Spec}\mathcal{O}_x\to Y$ 的概形）。$\mathfrak{g}^{*}$ 的正则元构成开子集 $\mathfrak{g}^{*}_{reg}$，使得 $\operatorname{codim}(\mathfrak{g}^{*}\setminus\mathfrak{g}^{*}_{reg})>1$。于是可把 $\operatorname{Paths}(\mathfrak{g}^{*})$ 换成 $\operatorname{Paths}(\mathfrak{g}^{*}_{reg})$。因态射 $\mathfrak{g}^{*}_{reg}\to\mathcal{C}$ 光滑且满，并且 $G$ 在其纤维上的作用传递，故得证。

\bdp{2.4.2} 由 1.2.2，$\mathfrak{z}_x^{cl}=P_{(\mathfrak{g}\otimes K_x,G(\mathcal{O}_x))}$ 是泊松代数。事实上，$\mathfrak{z}_x^{cl}$ 上的泊松括号为零，因为 1.2.7 中的态射 $\bar{a}^{cl}:(\overline{\operatorname{Sym}}(\mathfrak{g}\otimes K_x))^{\mathfrak{g}\otimes K_x}\to\mathfrak{z}_x^{cl}$ 是满射（例如可由 2.4.1 中 $\mathfrak{z}_x^{cl}$ 的描述推出），并且 $(\overline{\operatorname{Sym}}(\mathfrak{g}\otimes K_x))^{\mathfrak{g}\otimes K_x}$ 是 $\overline{\operatorname{Sym}}(\mathfrak{g}\otimes K_x)$ 的泊松中心。

\emph{注}（读者可跳过）。事实上，对任意代数群 $G$，自然态射 $\bar{a}^{cl}:(\overline{\operatorname{Sym}}(\mathfrak{g}\otimes K_x))^{G(K_x)}\to\mathfrak{z}_x^{cl}=\mathfrak{z}_{\mathfrak{g}}^{cl}(\mathcal{O}_x)$ 是满射，因此 $\mathfrak{z}_x^{cl}$ 上的泊松括号为零。下面的证明是 Feigin--Frenkel [FF92] 第 200--202 页论证的「经典极限」。把 $\mathcal{O}_x$ 与 $K_x$ 分别等同于 $\mathcal{O}:=\mathbb{C}[[t]]$ 与 $K:=\mathbb{C}((t))$。设 $f$ 是 $\mathfrak{g}^{*}\otimes\mathcal{O}$ 上的 $G(\mathcal{O})$-不变正则函数。须把它延拓为归纳概形 $\mathfrak{g}^{*}\otimes K:=\varinjlim\mathfrak{g}^{*}\otimes t^{-n}\mathcal{O}$ 上的 $G(K)$-不变正则函数 $\tilde{f}$（事实上 $\mathfrak{g}^{*}$ 可换成任意有限维 $G$-模）。至于

\SourcePage{23}

$\varphi\in\mathfrak{g}^{*}((t))$，定义 $h_{\varphi}\in\mathbb{C}((\zeta))$ 为
\[
h_{\varphi}(\zeta)=f\left(\sum_{k=0}^{N}\varphi^{(k)}(\zeta)t^{k}/k!\right),
\]
其中 $N$ 充分大（$h_{\varphi}$ 良定义，因为存在 $m$ 使得 $f$ 来自 $\mathfrak{g}^{*}\otimes(\mathcal{O}/t^{m}\mathcal{O})$ 上的函数）。把 $h_{\varphi}(\zeta)$ 写成 $\sum_n h_n(\varphi)\zeta^{n}$。函数 $h_n:\mathfrak{g}^{*}\otimes K\to\mathbb{C}$ 是 $G(K)$-不变的。置 $\tilde{f}:=h_0$。

\bdp{2.4.3} 由 2.3 与 1.2.6，有态射
\[
h_x^{cl}:\mathfrak{z}_x^{cl}\to\Gamma(\operatorname{Bun}_G,P),
\]
它类似于 1.2.3 中的态射 $h^{cl}$。为把它与 $h_X^{cl}$ 比较，考虑仿射概形的闭嵌入 $\operatorname{Hitch}(X)\hookrightarrow\operatorname{Hitch}_x$，它把 $\mathcal{C}_{\omega_X}$ 的整体截面映到其在 $x$ 的形式邻域上的限制。设 $\theta_x^{cl}:\mathfrak{z}_x^{cl}\to\mathfrak{z}^{cl}(X)$ 为相应的分次代数的满态射。容易看出
\[
h_x^{cl}=h_X^{cl}\theta_x^{cl}.
\]
因 $\mathfrak{z}_x^{cl}$ 上的泊松括号为零（见 2.4.2），且 $h_x^{cl}$ 是泊松代数态射，$\operatorname{Im}h_x^{cl}=\operatorname{Im}h_X^{cl}$ 上的泊松括号也为零。于是证明了 2.2.4(i)。

\bdsec{2.5. 量子化 I}

\bdp{2.5.1} 设 $\widetilde{\mathfrak{g}\otimes K_x}$ 为 $\mathfrak{g}\otimes K_x$ 的 Kac--Moody 中心扩张，中心为 $\mathbb{C}$，由余圈 $(u,v)\mapsto\operatorname{Res}_x c(du,v)$ 定义，$u,v\in\mathfrak{g}\otimes K_x$，其中
\begin{equation}\tag{18}
c(a,b):=-\frac{1}{2}\operatorname{Tr}(\operatorname{ad}_a\cdot\operatorname{ad}_b),\qquad a,b\in\mathfrak{g}.
\end{equation}
作为向量空间，$\widetilde{\mathfrak{g}\otimes K_x}$ 等于 $\mathfrak{g}\otimes K_x\oplus\mathbb{C}\cdot 1$。我们通过对 $g\in G(K_x)$ 指定如下自同构，来定义 $G(K_x)$ 在 $\widetilde{\mathfrak{g}\otimes K_x}$ 上的伴随作用\footnote{一旦有了李代数为 $\widetilde{\mathfrak{g}\otimes K_x}$ 的 $G(K_x)$ 的中心扩张，作用 (19) 就成为真正的伴随作用（$\widetilde{\mathfrak{g}\otimes K_x}$ 的一个在 $\mathbb{C}\cdot 1$ 与 $\mathfrak{g}\otimes K_x$ 上恒等的自同构必为恒等自同构，因为 $\operatorname{Hom}(\mathfrak{g}\otimes K_x,\mathbb{C})=0$）。}：
% [Tutor Note: 原书 p.23 末句的公式 (19) 起于 p.24，不在本片段范围内。]

\SourcePage{24}
$\tilde{\mathfrak{g}}\otimes K_x$ 的：

\begin{equation}\tag{19}
\mathbf{1}\mapsto\mathbf{1},\qquad
u\mapsto gug^{-1}+\operatorname{Res}_x c(u,g^{-1}dg)\cdot\mathbf{1},
\quad u\in\mathfrak{g}\otimes K_x .
\end{equation}

特别地，我们有 Harish-Chandra 对 $\bigl(\tilde{\mathfrak{g}}\otimes K_x, G(\mathcal{O}_x)\bigr)$。
它是 $\bigl(\mathfrak{g}\otimes K_x, G(\mathcal{O}_x)\bigr)$ 被 $\mathbb{C}$ 的中心扩张。记
\[
\mathfrak{z}_x=\mathfrak{z}_{\mathfrak{g}}(\mathcal{O}_x):=D'_{(\mathfrak{g}\otimes K_x,G(\mathcal{O}_x))},
\]
其中 $D'$ 的含义同 1.2.5。

\bdp{2.5.2}\emph{定理.}（[FF92]）。
\begin{enumerate}[label=(\roman*),leftmargin=2.2em]
\item 代数 $\mathfrak{z}_x$ 交换。
\item $\bigl(\tilde{\mathfrak{g}}\otimes K_x, G(\mathcal{O}_x)\bigr)$ 满足扭局部量子化条件（见 1.2.5）。
  也就是说，典范态射 $\sigma_{\mathfrak{z}_x}: \operatorname{gr}\mathfrak{z}_x\to\mathfrak{z}^{cl}_x$ 是同构。\hfill$\square$
\end{enumerate}

\emph{注.} 定理的论断 (i) 在 [FF92] 中对任意代数群 $G$ 有证明，
也对由 $\mathfrak{g}$ 上对称不变双线性型定义的 $\mathfrak{g}\otimes K_x$ 的任意中心扩张有证明。
此外，[FF92] 还证明：完备化扭泛包络代数 $\overline{U}'(\mathfrak{g}\otimes K_x)$ 的中心的
$\pi_0(G(K_x))$-不变部分满射到 $\mathfrak{z}_x$ 上。
Feigin–Frenkel 关于 (i) 的证明的一个版本将在 2.9.3–2.9.5 中给出。
我们已在 2.4.2 末尾的注中解释了其证明的「经典极限」。

\bdp{2.5.3} 线丛 $\omega_{\operatorname{Bun}_G}$ 在 $\operatorname{Bun}_{G,\underline{x}}$ 上定义一个 $G(\mathcal{O}_x)$-等变丛。
$\operatorname{Bun}_{G,\underline{x}}$ 上的 $(\mathfrak{g}\otimes K_x, G(\mathcal{O}_x))$-作用典范地提升为
该线丛上的 $(\tilde{\mathfrak{g}}\otimes K_x, G(\mathcal{O}_x))$-作用，使得 $\mathbf{1}$ 的作用为乘以 2。
事实上，由 2.1.1 有 $\omega_{\operatorname{Bun}_G}=f^*(\det R\Gamma)$，
其中 $f: \operatorname{Bun}_G\to\operatorname{Bun}_{SL(\mathfrak{g})}$ 由伴随表示 $G\to SL(\mathfrak{g})$ 诱导，
而 $\det R\Gamma$ 是 $\operatorname{Bun}_{SL(\mathfrak{g})}$ 上的行列式线丛。
另一方面，众所周知（例如见 [BLa94]），$\det R\Gamma$ 到 $\operatorname{Bun}_{SL_n,\underline{x}}$ 的拉回上
配有 $sl_n(K_x)$ 的 Kac–Moody 扩张（水平 $-1$）的作用。

\SourcePage{25}
\emph{注.} 事实上，该扩张的作用可积分为 $SL_n(K_x)$ 的某个中心扩张的作用（例如见 [BLa94]）。
因此得到典范的中心扩张
\begin{equation}\tag{20}
0\to\mathbb{G}_m\to\hat{G}(K_x)\to G(K_x)\to 0
\end{equation}
它在 $\omega_{\operatorname{Bun}_G}$ 到 $\operatorname{Bun}_{G,\underline{x}}$ 的拉回上的作用
使得 $\lambda\in\mathbb{G}_m$ 的作用为乘以 $\lambda$。
扩张 $0\to\mathbb{C}\to\tilde{\mathfrak{g}}\otimes K_x\to\mathfrak{g}\otimes K_x\to 0$
是对应于 (20) 的李代数扩张的一半。
在第 4 章中，我们将引入 $\omega_{\operatorname{Bun}_G}$ 的一个平方根\footnote{该平方根与扩张 (21) 依赖于 $\omega_X$ 的平方根的选取。}（Pfaffian 丛）以及中心扩张
\begin{equation}\tag{21}
0\to\mathbb{G}_m\to\tilde{G}(K_x)\to G(K_x)\to 0
\end{equation}
（见 4.4.8），它是 (20) 的一个平方根。
这些平方根对我们比 $\omega_{\operatorname{Bun}_G}$ 与 (20) 更重要，因此我们不给出 $\tilde{G}(K_x)$ 的精确定义。

\bdp{2.5.4} 由 2.5.3 与 1.2.5，我们有滤代数的一个典范态射
\[
h_x: \mathfrak{z}_x\to\Gamma(\operatorname{Bun}_G,D').
\]
在 2.7.5 中我们将证明下述定理。

\bdp{2.5.5}\emph{定理.} 我们的数据满足扭整体量子化条件（见 1.2.5）。\hfill$\square$

如 1.2.3 所解释，由于局部与整体量子化条件都满足，
我们得到 2.2.5 的量子化问题的一个解 $\mathfrak{z}^{(x)}(X)$：
置 $\mathfrak{z}^{(x)}(X)=h_x(\mathfrak{z}_x)$，并给 $\mathfrak{z}^{(x)}(X)$ 配上
由 $\Gamma(\operatorname{Bun}_G,D')$ 上的滤链诱导的滤链
（2.5.5 意味着它也可由 $\mathfrak{z}_x$ 上的滤链诱导）；
于是象征映射把 $\operatorname{gr}\mathfrak{z}^{(x)}(X)$ 与 $h_x^{cl}(\mathfrak{z}^{cl}_x)$ 等同，
且由 2.4.3 有
$h_x^{cl}(\mathfrak{z}^{cl}_x)=h_X^{cl}(\mathfrak{z}^{cl}_x(X))\simeq\mathfrak{z}^{cl}(X)$。

\SourcePage{26}
定理 2.5.5 的证明基于 2.2.5 中量子化问题的解的第二种构造；
它还表明 $\mathfrak{z}^{(x)}(X)$ 不依赖于 $x$。

\emph{注.} 若 $G$ 单连通，则 2.5.5 立即由 2.2.4(v) 得到。

\bdsec{2.6. \texorpdfstring{$\mathcal{D}_X$}{D\_X}-概形的一般性讨论}

\bdp{2.6.1} 设 $X$ 为任意光滑连通代数簇。
一个 $\mathcal{D}_X$-概形是指带有沿 $X$ 的平坦联络的 $X$-概形。
在 $X$ 上仿射的 $\mathcal{D}_X$-概形是交换 $\mathcal{D}_X$-代数的谱
（交换 $\mathcal{D}_X$-代数即带有平坦联络的拟凝聚 $\mathcal{O}_X$-代数）。
$\mathcal{O}_X$-代数 $\mathcal{A}$ 在 $x\in X$ 处的纤维记作 $\mathcal{A}_x$；
这尤其适用于 $\mathcal{D}_X$-代数。
对 $\mathbb{C}$-代数 $C$，记 $C_X$ 为相应的「常值」 $\mathcal{D}_X$-代数
（即 $C_X$ 是 $C\otimes\mathcal{O}_X$ 配以显然的联络）。

\bdp{2.6.2}\emph{命题.} 设 $X$ 完备。
\begin{enumerate}[label=(\roman*),leftmargin=2.2em]
\item 函子 $C\rightsquigarrow C_X$ 有左伴随函子：对 $\mathcal{D}_X$-代数 $\mathcal{A}$，
  存在 $\mathbb{C}$-代数 $H_\nabla(X,\mathcal{A})$，使得
  \begin{equation}\tag{22}
  \operatorname{Hom}(\mathcal{A},C_X)=\operatorname{Hom}(H_\nabla(X,\mathcal{A}),C)
  \end{equation}
  对任意 $\mathbb{C}$-代数 $C$ 成立。
\item 典范投影 $\theta_{\mathcal{A}}: \mathcal{A}\to H_\nabla(X,\mathcal{A})_X$ 满。
  于是 $H_\nabla(X,\mathcal{A})_X$ 是 $\mathcal{A}$ 的最大「常值」商 $\mathcal{D}_X$-代数。
  特别地，对任意 $x\in X$，态射
  $\theta_{\mathcal{A}_x}: \mathcal{A}_x\to(H_\nabla(X,\mathcal{A})_X)_x=H_\nabla(X,\mathcal{A})$ 满。
\end{enumerate}

\emph{注.} (i) 这里的代数不假定交换、结合等。我们将在 $\mathcal{A}$ 交换的情形下需要该命题。

(ii) 设 $\mathcal{A}$ 交换（这是「交换结合含幺」的缩写）。
则按命题的论断 (ii)，$H_\nabla(X,\mathcal{A})$ 交换。
若 $C$ 也假定交换，则 (22) 只是说
$\operatorname{Spec} H_\nabla(X,\mathcal{A})$ 是 $\operatorname{Spec}\mathcal{A}$ 的水平截面概形。

\SourcePage{27}
从几何观点看，这样的概形存在且仿射是清楚的：
$\operatorname{Spec}\mathcal{A}$ 的所有截面构成仿射概形 $S$
（这里用到 $X$ 的完备性；否则 $S$ 将是归纳概形，见下一个注），
而水平截面构成 $S$ 的闭子概形。
$\theta_{\mathcal{A}_x}$ 与 $\theta_{\mathcal{A}}$ 的满性意味着态射
$\operatorname{Spec} H_\nabla(X,\mathcal{A})\to\operatorname{Spec}\mathcal{A}_x$
与 $X\times\operatorname{Spec} H_\nabla(X,\mathcal{A})\to\operatorname{Spec}\mathcal{A}$ 是闭浸入。

(iii) 若 $X$ 任意（不必完备），则由 (22) 定义的 $H_\nabla(X,\mathcal{A})$
可表为代数关于某个有向满射族的投射极限。
于是若 $\mathcal{A}$ 交换，则 $\operatorname{Spec}\mathcal{A}$ 的水平截面空间是归纳仿射归纳概形\footnote{从几何观点看这也很清楚。事实上，水平截面构成所有截面空间 $S_X$ 中的闭子空间。若 $X$ 仿射，$S_X$ 当然是归纳概形。一般情形下 $X$ 可被开仿射子概形 $U_1,\dots,U_n$ 覆盖；于是 $S_X$ 是诸 $S_{U_i}$ 之积的闭子空间。}。

\emph{证明.} (a) 记 $\mathcal{M}(X)$ 为 $\mathcal{D}_X$-模的范畴，
$\mathcal{M}_{const}(X)$ 为常值 $\mathcal{D}_X$-模的满子范畴，
即同构于 $V\otimes\mathcal{O}_X$（$V$ 为某向量空间）的 $\mathcal{D}_X$-模
（实际上函子 $V\mapsto V\otimes\mathcal{O}_X$ 是向量空间范畴与 $\mathcal{M}_{const}(X)$ 之间的等价）。
我们断言嵌入 $\mathcal{M}_{const}(X)\to\mathcal{M}(X)$ 有左伴随函子，
即对 $\mathcal{F}\in\mathcal{M}(X)$ 存在 $\mathcal{F}_\nabla\in\mathcal{M}_{const}(X)$，使得
$\operatorname{Hom}(\mathcal{F},\mathcal{E})=\operatorname{Hom}(\mathcal{F}_\nabla,\mathcal{E})$
对 $\mathcal{E}\in\mathcal{M}_{const}(X)$ 成立。
只需对凝聚的 $\mathcal{F}$ 构造 $\mathcal{F}_\nabla$。
此时 $\mathcal{F}_\nabla:=(\operatorname{Hom}_{\mathcal{D}_X}(\mathcal{F},\mathcal{O}_X))^*\otimes\mathcal{O}_X$
（这里用到 $\dim\operatorname{Hom}_{\mathcal{D}_X}(\mathcal{F},\mathcal{O}_X)<\infty$，因为 $X$ 完备）。

(b) 由于 $\mathcal{O}_X$ 是不可约 $\mathcal{D}_X$-模，常值 $\mathcal{D}_X$-模的 $\mathcal{D}_X$-子模是常值的。
于是自然态射 $\mathcal{F}\to\mathcal{F}_\nabla$ 满。

(c) 若 $\mathcal{A}$ 是 $\mathcal{D}_X$-代数，$\mathcal{I}$ 是 $\mathcal{A}$ 中由
$\operatorname{Ker}(\mathcal{A}\to\mathcal{A}_\nabla)$ 生成的理想，
则 $\mathcal{A}/\mathcal{I}$ 是常值 $\mathcal{D}_X$-模 $\mathcal{A}_\nabla$ 的商。
所以 $\mathcal{A}/\mathcal{I}$ 常值，即 $\mathcal{A}/\mathcal{I}=H_\nabla(X,\mathcal{A})\otimes\mathcal{O}_X$，
其中 $H_\nabla(X,\mathcal{A})$ 为某向量空间。
$\mathcal{A}/\mathcal{I}$ 是 $\mathcal{D}_X$-代数，故 $H_\nabla(X,\mathcal{A})$ 是代数。
显然它满足 (22)。\hfill$\square$

\SourcePage{28}
\emph{注.} 偏几何的读者可把上面的注 (ii) 看作交换代数情形下命题的证明。
然而在 2.7.4 中，我们将在 $\mathcal{A}$ 交换而 $C=\Gamma(\operatorname{Bun}_G,D')$ 并不显然交换的情形下应用 (22)。
此时只需注意：态射 $\mathcal{A}\to C\otimes\mathcal{O}_X$ 的像形如 $C'\otimes\mathcal{O}_X$
（见命题证明的 (b) 部分），且 $C'$ 交换，因为 $\mathcal{A}$ 交换。
也可把 (22) 用于 $C:=\Gamma(\operatorname{Bun}_G,D')$ 中由所有 $x\in X$ 的态射
$h_x: \mathfrak{z}_x\to\Gamma(\operatorname{Bun}_G,D')$ 的像生成的子代数
（这个 $C$ 「显然」交换；见 2.9.1）。
事实上，用 2.2.4(v) 可以证明 $\Gamma(\operatorname{Bun}_G,D')$ 交换
（由 2.2.4(v) 可推出：对任意连通分支
$\operatorname{Bun}_G^\gamma\subset\operatorname{Bun}_G$ 与任意 $x\in X$，
由 $h_x$ 诱导的态射 $\mathfrak{z}_x\to\Gamma(\operatorname{Bun}_G^\gamma,D')$ 满）。

\bdp{2.6.3} 本小节中所有代数均假定交换。
遗忘函子 $\{\mathcal{D}_X\text{-代数}\}\to\{\mathcal{O}_X\text{-代数}\}$ 有显然的左伴随函子 $\mathcal{J}$
（$\mathcal{J}\mathcal{A}$ 是由 $\mathcal{O}_X$-代数 $\mathcal{A}$ 生成的 $\mathcal{D}_X$-代数）。
我们断言 $\operatorname{Spec}\mathcal{J}\mathcal{A}$ 正是
\emph{$\operatorname{Spec}\mathcal{A}$ 的截面的 $\infty$-射流概形}。
特别地，这意味着 $\operatorname{Spec}(\mathcal{J}\mathcal{A})_x$ 的 $\mathbb{C}$-点
与截面 $\operatorname{Spec}\mathcal{O}_x\to\operatorname{Spec}\mathcal{A}$ 之间有典范的一一对应
（其中 $\mathcal{O}_x$ 是 $x$ 处局部环的形式完备化）。
更确切地说，我们要构造函子性的双射
\begin{equation}\tag{23}
\operatorname{Hom}_{\mathcal{O}_X}(\mathcal{J}\mathcal{A},\mathcal{B})
\overset{\sim}{\longrightarrow}
\operatorname{Hom}_{\mathcal{O}_X}(\mathcal{A},\hat{\mathcal{B}})
\end{equation}
其中 $\mathcal{B}$ 是（拟凝聚）$\mathcal{O}_X$-代数，$\hat{\mathcal{B}}$ 是
$\mathcal{O}_X\otimes_{\mathbb{C}}\mathcal{B}$ 关于理想
$\operatorname{Ker}(\mathcal{O}_X\otimes_{\mathbb{C}}\mathcal{B}\to\mathcal{B})$ 的完备化。
这里 $\hat{\mathcal{B}}$ 配有由态射
$\mathcal{O}_X\to\mathcal{O}_X\otimes_{\mathbb{C}}\mathcal{B}$（$a\mapsto a\otimes 1$）给出的
$\mathcal{O}_X$-代数结构。
让我们暂时去掉 $\mathcal{D}_X$-代数定义中的拟凝聚假定。
则 $\hat{\mathcal{B}}$ 是 $\mathcal{D}_X$-代数
（$\hat{\mathcal{B}}$ 上的联络来自 $\mathcal{O}_X\otimes_{\mathbb{C}}\mathcal{B}$ 上的联络，
该联络使 $1\otimes\mathcal{B}$ 的截面水平）。
于是 $\operatorname{Hom}_{\mathcal{O}_X}(\mathcal{A},\hat{\mathcal{B}})
=\operatorname{Hom}_{\mathcal{D}_X}(\mathcal{J}\mathcal{A},\hat{\mathcal{B}})$，
而要构造 (23)，只需构造函子性的双射
\SourcePage{29}
\begin{equation}\tag{24}
\operatorname{Hom}_{\mathcal{O}_X}(\mathcal{R},\mathcal{B})
\longleftarrow
\operatorname{Hom}_{\mathcal{D}_X}(\mathcal{R},\hat{\mathcal{B}})
\end{equation}
对任意 $\mathcal{D}_X$-代数 $\mathcal{R}$ 与 $\mathcal{O}_X$-代数 $\mathcal{B}$ 成立
（即证明函子 $\mathcal{B}\mapsto\hat{\mathcal{B}}$ 是遗忘函子
$\{\mathcal{D}_X\text{-代数}\}\to\{\mathcal{O}_X\text{-代数}\}$ 的\emph{右}伴随）。
映射 (24) 来自显然的态射 $\hat{\mathcal{B}}\to\mathcal{B}$。
读者容易证明 (24) 是双射。

对 $\mathcal{D}_X$-代数 $\mathcal{A}$ 与 $\mathbb{C}$-代数 $C$，有
\[
\operatorname{Hom}_{\mathcal{D}_X\text{-alg}}(\mathcal{J}\mathcal{A},C\otimes\mathcal{O}_X)
=\operatorname{Hom}_{\mathcal{O}_X\text{-alg}}(\mathcal{A},C\otimes\mathcal{O}_X).
\]
这意味着：\emph{典范态射 $\operatorname{Spec}\mathcal{J}\mathcal{A}\to\operatorname{Spec}\mathcal{A}$
把 $\operatorname{Spec}\mathcal{J}\mathcal{A}$ 的水平截面归纳概形
与 $\operatorname{Spec}\mathcal{A}$ 的（所有）截面的归纳概形等同}。
若 $X$ 完备，则由 2.6.2，这些空间实际上是概形。

最后我们指出，本小节的结果可按显然的方式整体化。
遗忘函子 $\{\mathcal{D}_X\text{-概形}\}\to\{X\text{-概形}\}$ 有右伴随函子
$\mathcal{J}: \{X\text{-概形}\}\to\{\mathcal{D}_X\text{-概形}\}$。
对 $X$-概形 $Y$，$\mathcal{J}Y$ 是 $Y$ 的截面的 $\infty$-射流概形。
对 $\mathcal{O}_X$-代数 $\mathcal{A}$，有 $\mathcal{J}\operatorname{Spec}\mathcal{A}=\operatorname{Spec}\mathcal{J}\mathcal{A}$。
典范态射 $\mathcal{J}Y\to Y$ 把 $\mathcal{J}Y$ 的水平截面空间\footnote{最一般的情形下，「空间」指「函子 $\{\mathbb{C}\text{-代数}\}\to\{\text{集合}\}$」。}
与 $Y$ 的（所有）截面空间等同。
若 $X$ 完备且 $Y$ 拟射影，则我们的空间是概形。

\bdp{2.6.4} 设 $(\mathfrak{l},P)$ 是 1.2.6 意义下的 Harish-Chandra 对
（故 $P$ 可为任意仿射群概形；我们不假定它有限型\footnote{由下面的定义可见，$\operatorname{Lie} P$ 在 $\mathfrak{l}$ 中的余维有限（等于 $\dim X$）。}）。

\emph{定义.} $X$ 上的 $(\mathfrak{l},P)$-结构是指态射 $\pi: X^{\wedge}\to X$
连同 $(\mathfrak{l},P)$ 在 $X^{\wedge}$ 上的作用，使得
\begin{enumerate}[label=(\roman*),leftmargin=2.2em]
\item $X^{\wedge}$ 是 $X$ 上的 $P$-torsor。
\item $\mathfrak{l}$ 在 $X^{\wedge}$ 上的作用形式自由且传递，即给出同构
  $\mathfrak{l}\hat{\otimes}\mathcal{O}_{X^{\wedge}}\overset{\sim}{\to}\Theta_{X^{\wedge}}$。
\end{enumerate}

\SourcePage{30}
\emph{注.} 设 $L$ 是李代数为 $\mathfrak{l}$、$L_{red}=P$ 的群归纳概形（见 7.11.2(v)）。
% [Tutor Note: 原书 p.30 此处拼作「homogenuous」，疑为「homogeneous」之笔误，底本如此。]
考虑齐性空间 $P\backslash L=\operatorname{Spf}\mathcal{O}$，
其中 $\mathcal{O}=\mathcal{O}_{(\mathfrak{l},P)}=(U\mathfrak{l}/(U\mathfrak{l})\mathfrak{p})^*$。
取 $x\in X$ 并选 $x^{\wedge}\in\pi^{-1}(x)$。
映射 $L\to X^{\wedge}$，$l\mapsto lx^{\wedge}$，给出态射
$\alpha_{x^{\wedge}}: \operatorname{Spf}\mathcal{O}\to X$，
它把 $\operatorname{Spf}\mathcal{O}$ 与 $x$ 的形式邻域等同。
对 $l\in L$，$a\in\operatorname{Spf}\mathcal{O}$，有 $\alpha_{lx^{\wedge}}(a)=\alpha_{x^{\wedge}}(al)$。
注意：若 $P$ 在 $\mathcal{O}$ 上的作用忠实，则 $x^{\wedge}$ 由 $\alpha_{x^{\wedge}}$ 唯一确定。

\bdp{2.6.5}\emph{例.} 置 $\mathcal{O}=\mathcal{O}_n:=\mathbb{C}[[t_1,..,t_n]]$。
$\mathbb{C}$-代数 $\mathcal{O}$ 的自同构群自然地是 $\mathbb{C}$ 上一个仿射群概形
$\operatorname{Aut}^0\mathcal{O}$ 的 $\mathbb{C}$-点群。
记 $\operatorname{Aut}\mathcal{O}$ 为这样的群归纳概形：
对任意 $\mathbb{C}$-代数 $R$，$(\operatorname{Aut}\mathcal{O})(R)$ 是拓扑 $R$-代数
$R\hat{\otimes}\mathcal{O}=R[[t_1,..,t_n]]$ 的自同构群。
于是 $\operatorname{Aut}^0\mathcal{O}$ 是 $\operatorname{Aut}\mathcal{O}$ 的群子概形；
事实上 $\operatorname{Aut}^0\mathcal{O}=(\operatorname{Aut}\mathcal{O})_{red}$。
有 $\operatorname{Lie}\operatorname{Aut}\mathcal{O}=\operatorname{Der}\mathcal{O}$，
$\operatorname{Lie}\operatorname{Aut}^0\mathcal{O}=\operatorname{Der}^0\mathcal{O}
:=\mathfrak{m}_{\mathcal{O}}\cdot\operatorname{Der}\mathcal{O}$。
因此 $\operatorname{Aut}\mathcal{O}$ 是对应于 Harish-Chandra 对
$\operatorname{Aut}^{HC}\mathcal{O}:=(\operatorname{Der}\mathcal{O},\operatorname{Aut}^0\mathcal{O})$ 的群归纳概形。
我们滥用记号，将写 $\operatorname{Aut}\mathcal{O}$ 而不用 $\operatorname{Aut}^{HC}\mathcal{O}$。

正如 Gelfand 与 Kazhdan 所解释（见 [GK]、[GKF] 与 [BR]），
任意 $n$ 维光滑代数簇 $X$ 带有典范的\footnote{事实上，$X$ 上的 $\operatorname{Aut}\mathcal{O}$-结构在同构意义下唯一（同构也唯一，这由 2.6.4 的注得出）。}
$\operatorname{Aut}\mathcal{O}$-结构。
空间 $X^{\wedge}=X^{\wedge}_{can}$ 是 $X$ 上「形式坐标系」的空间。
换言之，$X^{\wedge}$ 的一个 $\mathbb{C}$-点是微分非零的态射 $\operatorname{Spec}\mathcal{O}\to X$，
而 $X^{\wedge}$ 的一个 $R$-点是 $R$-态射
$\alpha: \operatorname{Spec}(R\hat{\otimes}\mathcal{O})\to X\otimes R$，
其微分在 $\operatorname{Spec} R$ 的每一点上非零。
群归纳概形 $\operatorname{Aut}\mathcal{O}$ 以显然的方式作用在 $X^{\wedge}$ 上，
且有投影 $\pi: X^{\wedge}\to X$，$\alpha\mapsto\alpha(0)$。
容易看出 $X^{\wedge}$（连同上列结构）是 $X$ 上的 $\operatorname{Aut}\mathcal{O}$-结构。

我们将用到 $n=1$ 时的典范 $\operatorname{Aut}\mathcal{O}_n$-结构，
即 $X$ 为曲线时，此时 $\mathcal{O}=\mathbb{C}[[t]]$。
这里群 $\operatorname{Aut}\mathcal{O}$ 如下。
有满同态
$\operatorname{Aut}^0\mathcal{O}\to\operatorname{Aut}(t\mathcal{O}/t^2\mathcal{O})=\mathbb{G}_m$，
\SourcePage{31}
我们称之为 $\operatorname{Aut}^0\mathcal{O}$ 的\emph{标准特征标}；其核是投射幂幺的。
对 $\mathbb{C}$-代数 $R$，$R[[t]]$ 的自同构由 $t\mapsto\sum_i c_i t^i$ 定义，
其中 $c_1\in R^*$，$c_0$ 幂零。
于是 $\operatorname{Aut}\mathcal{O}$ 是概形
$\operatorname{Spec}\mathbb{C}[c_0,c_1,c_1^{-1},c_2,c_3,\dots]/(c_0^k)$（$k\in\mathbb{N}$）的并。
$\operatorname{Aut}^0\mathcal{O}$ 是由 $c_0=0$ 定义的 $\operatorname{Aut}\mathcal{O}$ 的群子概形。

另一些 $(\mathfrak{l},P)$-结构的例子可在 ?? 中找到。
% [Tutor Note: 原书 p.31 此处引用印作「??」，底本如此，保留原样。]

\bdp{2.6.6} 设 $X$ 是带有 $(\mathfrak{l},P)$-结构 $X^{\wedge}$ 的代数簇
（我们将在下述情形下应用这些构造：$X$ 为曲线，$\mathfrak{l}=\operatorname{Der}\mathcal{O}$，
$P=\operatorname{Aut}^0\mathcal{O}$（或 $\operatorname{Aut}^0\mathcal{O}$ 的某个覆叠），
$\mathcal{O}:=\mathbb{C}[[t]]$）。
记 $\mathcal{M}(X,\mathcal{O})$ 为 $X$ 上 $\mathcal{O}$-模的范畴，
$\mathcal{M}^l(X)$ 为左 $\mathcal{D}$-模的范畴。
对 $\mathcal{F}_X\in\mathcal{M}(X,\mathcal{O})$，其到 $X^{\wedge}$ 的拉回
$\mathcal{F}_{X^{\wedge}}$ 是 $X^{\wedge}$ 上的 $P$-等变 $\mathcal{O}$-模。
若 $\mathcal{F}_X$ 实际上是左 $\mathcal{D}_X$-模，则 $\mathcal{F}_{X^{\wedge}}$ 还是 $\mathfrak{l}$-等变的
（因为由 2.6.4(ii)，$\mathcal{O}_{X^{\wedge}}$-模上的 $\mathfrak{l}$-作用等同于平坦联络）。
函子
$\mathcal{M}(X,\mathcal{O})\to\{X^{\wedge}\text{ 上的 }P\text{-等变 }\mathcal{O}\text{-模}\}$、
$\mathcal{M}^l(X)\to\{X^{\wedge}\text{ 上的 }(\mathfrak{l},P)\text{-等变 }\mathcal{O}\text{-模}\}$
是张量范畴的等价。

有忠实正合张量函子
\begin{equation}\tag{25}
\mathcal{M}(P)\longrightarrow\mathcal{M}(X,\mathcal{O}),\qquad
\mathcal{M}(\mathfrak{l},P)\longrightarrow\mathcal{M}^l(X)
\end{equation}
它们把表示 $V$ 送到 $\mathcal{O}_X$-模或 $\mathcal{D}_X$-模 $V_X$，
使得 $V_{X^{\wedge}}$ 等于 $V\otimes\mathcal{O}_{X^{\wedge}}$
（$P$-模或 $(\mathfrak{l},P)$-模的张量积）。
换言之，$\mathcal{O}_X$-模 $V_X$ 是 $V$ 用 $P$-torsor $X^{\wedge}$ 扭出的结果。
因此任意带 $P$-作用的代数 $A$ 给出 $\mathcal{O}_X$-代数 $A_X$；
若 $A$ 实际上带有 $(\mathfrak{l},P)$-作用，则 $A_X$ 是 $\mathcal{D}_X$-代数。
类似地，任意带 $P$-作用的概形 $H$（简称\emph{$P$-概形}）给出 $X$-概形 $H_X$。
若 $H$ 实际上是 $(\mathfrak{l},P)$-概形，则 $H_X$ 是 $\mathcal{D}_X$-概形。
有 $(\operatorname{Spec} A)_X=\operatorname{Spec}(A_X)$。

\emph{注.} (i) 函子 $\mathcal{M}(\mathfrak{l},P)\longrightarrow\mathcal{M}^l(X)$
与 $X^{\wedge}$ 上 $(\mathfrak{l},P)$-作用的局部化函子 $\Delta$ 一致（见 1.2.4）。

(ii) 函子 (25) 有右伴随，它把 $\mathcal{O}_X$-模或 $\mathcal{D}_X$-模 $\mathcal{F}_X$ 送到向量空间
\SourcePage{32}
$\Gamma(X^{\wedge},\mathcal{F}_{X^{\wedge}})$，其上配有显然的 $P$-模或 $(\mathfrak{l},P)$-模结构。
若考虑代数而非模，同样的伴随性成立。

(iii) 设 $C$ 是 $\mathbb{C}$-代数；把 $C$ 看作带平凡 $\operatorname{Aut}\mathcal{O}$-作用的 $(\mathfrak{l},P)$-代数。
则 $C_X$ 是 2.6.1 中的「常值」 $\mathcal{D}_X$-代数。

\bdp{2.6.7} 遗忘函子 $\{(\mathfrak{l},P)\text{-代数}\}\to\{P\text{-代数}\}$
有左伴随（诱导）函子 $\mathcal{J}$。
对 $P$-代数 $A$ 有典范同构
\begin{equation}\tag{26}
(\mathcal{J}A)_X=\mathcal{J}(A_X).
\end{equation}
事实上，自然的 $\mathcal{O}_X$-代数态射 $A_X\to(\mathcal{J}A)_X$ 诱导
$\mathcal{D}_X$-代数态射 $\mathcal{J}(A_X)\to(\mathcal{J}A)_X$。
为证它是同构，用 $\mathcal{J}$ 与 $A\mapsto A_X$ 的伴随性（见 2.6.3 与 2.6.6 的注 (ii)）。

上面陈述的几何版本如下。
遗忘函子 $\{(\mathfrak{l},P)\text{-概形}\}\to\{P\text{-概形}\}$ 有右伴随函子\footnote{对仿射概形，这只是上面关于 $P$-代数的陈述的改写。一般情形不能立即归结为仿射情形（$P$-概形未必有 $P$-不变仿射子概形的覆盖），但仿射情形已够我们使用。} $\mathcal{J}$。
对 $P$-代数 $A$ 有 $\mathcal{J}(\operatorname{Spec} A)=\operatorname{Spec}\mathcal{J}(A)$。
对任意 $P$-概形 $H$ 有 $(\mathcal{J}H)_X=\mathcal{J}(H_X)$。

\bdsec{2.7. 量子化 II}
从现在起 $\mathcal{O}:=\mathbb{C}[[t]]$，$K:=\mathbb{C}((t))$。

\bdp{2.7.1} 先考虑「经典」图景。
概形 $\operatorname{Hitch}_x$（$x\in X$）是在 $X$ 上仿射的 $\mathcal{D}_X$-概形
$\operatorname{Hitch}=\mathcal{J}C_{\omega_X}$ 的纤维；
记 $\mathfrak{z}^{cl}$ 为相应的 $\mathcal{D}_X$-代数。
由 2.6.3，投影 $\operatorname{Hitch}\to C_{\omega_X}$ 把 $\operatorname{Hitch}$ 的水平截面概形
与 $\operatorname{Hitch}(X)$ 等同。换言之
\[
\mathfrak{z}^{cl}(X)=H_\nabla(X,\mathfrak{z}^{cl}),
\]
而 2.4.3 中的投影 $\theta_x^{cl}: \mathfrak{z}_x^{cl}\to\mathfrak{z}^{cl}(X)$
正是命题 2.6.2(ii) 中的典范态射 $\theta_{\mathfrak{z}^{cl}_x}$。

\SourcePage{33}
把 $C$ 经由标准特征标 $\operatorname{Aut}^0\mathcal{O}\to\mathbb{G}_m$
看作 $\operatorname{Aut}^0\mathcal{O}$-概形（见 2.6.5）。
$X$-概形 $C_{\omega_X}$ 与 $C$ 的 $X^{\wedge}$-扭重合。
因此同构 (26) 诱导典范同构
\[
\mathfrak{z}^{cl}=\mathfrak{z}^{cl}_{\mathfrak{g}}(\mathcal{O})_X,
\]
其中 $\mathfrak{z}^{cl}_{\mathfrak{g}}(\mathcal{O})$ 是 $\operatorname{Aut}\mathcal{O}$-代数
$\mathcal{J}(\operatorname{Sym}\mathfrak{g})^G$，
而 $(\operatorname{Sym}\mathfrak{g})^G$ 上的 $\operatorname{Aut}^0\mathcal{O}$-作用
来自与 $(\operatorname{Sym}\mathfrak{g})^G$ 的分次所诱导的作用相反的 $\mathbb{G}_m$-作用（参看 2.2.2 的末尾）。

\bdp{2.7.2} 让我们转向「量子」图景。
置 $\mathfrak{z}_{\mathfrak{g}}(\mathcal{O}):=D'_{(\mathfrak{g}\otimes K,G(\mathcal{O}))}$。
这是交换代数（见 2.5.2(i)）。
$\operatorname{Aut}\mathcal{O}$ 作用在 $\mathfrak{z}_{\mathfrak{g}}(\mathcal{O})$ 上，
因为 $\mathfrak{z}_{\mathfrak{g}}(\mathcal{O})$ 是扭真空模 $\mathit{Vac}'$ 的自同态代数（见 1.2.5），
而 $\operatorname{Aut}\mathcal{O}$ 作用在 $\mathit{Vac}'$ 上。
（后一作用由两条性质刻画：它与 $\operatorname{Aut}\mathcal{O}$ 在 $\tilde{\mathfrak{g}}\otimes K$ 上的自然作用相容，
且真空向量不变；$\operatorname{Aut}\mathcal{O}$ 在 $\tilde{\mathfrak{g}}\otimes K$ 上的作用按拓扑意义理解，
即 $\operatorname{Aut}(\mathcal{O}\hat{\otimes}R)$ 作用在 $\tilde{\mathfrak{g}}\otimes K\hat{\otimes}R$ 上，
对任意交换 $\mathbb{C}$-代数 $R$。）考虑 $\mathcal{D}_X$-代数
\[
\mathfrak{z}=\mathfrak{z}_{\mathfrak{g}}:=\mathfrak{z}_{\mathfrak{g}}(\mathcal{O})_X,
\]
它对应于交换 $(\operatorname{Aut}\mathcal{O})$-代数 $\mathfrak{z}_{\mathfrak{g}}(\mathcal{O})$（见 2.6.5、2.6.6）。
其纤维是 2.5.1 中的代数 $\mathfrak{z}_x$。
标准论证表明，当 $x\in X$ 变动时，2.5.4 中的态射 $h_x$ 定义了
$\mathcal{O}_X$-代数态射 $h: \mathfrak{z}\to\Gamma(\operatorname{Bun}_G,D')_X$。

\bdp{2.7.3}\emph{水平性定理.} $h$ 水平，即它是 $\mathcal{D}_X$-代数态射。

证明见 2.8。

\bdp{2.7.4} 置
\begin{equation}\tag{27}
\mathfrak{z}(X)=\mathfrak{z}_{\mathfrak{g}}(X):=H_\nabla(X,\mathfrak{z}).
\end{equation}

\SourcePage{34}

按 2.6.2(i)，$\mathcal{D}_X$-代数态射 $h$ 诱导 $\mathbb{C}$-代数态射
\[
  h_X:\mathfrak{z}(X)\longrightarrow\Gamma(\operatorname{Bun}_G,D')
\]

我们将证明 $(\mathfrak{z}(X),h_X)$ 是 2.2.5 中量子化问题的一个解。
在此之前，须先定义 $\mathfrak{z}(X)$ 上的滤链，以及同构
$\sigma_{\mathfrak{z}(X)}:\operatorname{gr}\mathfrak{z}(X)\xrightarrow{\sim}\mathfrak{z}^{cl}(X)$。

$\mathfrak{z}_{\mathfrak{g}}(\mathcal{O})$ 上的典范滤链是
$\operatorname{Aut}\mathcal{O}$-不变的。同构
$\sigma_{\mathfrak{z}(\mathcal{O})}:\operatorname{gr}\mathfrak{z}_{\mathfrak{g}}(\mathcal{O})
\xrightarrow{\sim}\mathfrak{z}^{cl}_{\mathfrak{g}}(\mathcal{O})$（见 2.5.2(ii)）
与 $\operatorname{Aut}\mathcal{O}$-作用相容。因此 $\mathfrak{z}$ 带有水平滤链，
并且有 $\mathcal{D}_X$-代数的同构
\[
  \sigma_{\mathfrak{z}}:\operatorname{gr}\mathfrak{z}\xrightarrow{\sim}\mathfrak{z}^{cl}
\]
它在每条纤维上化为 2.5.2(ii) 的同构 $\sigma_{\mathfrak{z}_x}$。

通过 $\theta_{\mathfrak{z}}:\mathfrak{z}\to H_\nabla(X,\mathfrak{z})_X=\mathfrak{z}(X)_X$，
该滤链的像给出 $\mathfrak{z}(X)_X$ 上的水平滤链，
这等同于 $\mathfrak{z}(X)$ 上的滤链。考虑分次 $\mathcal{D}_X$-代数的满态射
$(\operatorname{gr}\theta_{\mathfrak{z}})\sigma_{\mathfrak{z}}^{-1}:\mathfrak{z}^{cl}\to
\operatorname{gr}\mathfrak{z}(X)_X$。由伴随（见 (22)），
它给出分次 $\mathbb{C}$-代数的满态射
\[
  j:\mathfrak{z}^{cl}(X)=H_\nabla(X,\mathfrak{z}^{cl})\longrightarrow\operatorname{gr}\mathfrak{z}(X).
\]

注意 $h_X$ 与滤链相容，且有交换图
\begin{equation}\tag{28}
\begin{tikzcd}[column sep=3.4em, row sep=2.4em]
\mathfrak{z}^{cl}(X) \arrow[r, hook, "h^{cl}_X"] \arrow[d, "j"]
  & \Gamma(\operatorname{Bun}_G,P) \arrow[d, "\sigma_{\operatorname{Bun}_G}"] \\
\operatorname{gr}\mathfrak{z}(X) \arrow[r, hook, "\operatorname{gr}h_X"]
  & \operatorname{gr}\Gamma(\operatorname{Bun}_G,D')
\end{tikzcd}
\end{equation}

因此 $j$ 是同构，$\operatorname{gr}h_X$（从而 $h_X$）是单射。
定义 $\sigma_{\mathfrak{z}(X)}:\operatorname{gr}\mathfrak{z}(X)\xrightarrow{\sim}
\mathfrak{z}^{cl}(X)$ 为 $\sigma_{\mathfrak{z}(X)}:=j^{-1}$。
三元组 $(\mathfrak{z}(X),h_X,\sigma_{\mathfrak{z}(X)})$ 是 2.2.5 中量子化问题的一个解。

\bdp{2.7.5} 让我们证明定理 2.5.5，并把 2.5.5 中的 $\mathfrak{z}^{(x)}(X)$
与 $\mathfrak{z}(X)$ 作比较。显然 $h_x=h_X\cdot\theta_{\mathfrak{z}_x}$，
其中 $\theta_{\mathfrak{z}_x}:\mathfrak{z}_x\to\mathfrak{z}(X)$ 由命题 2.6.2(ii) 定义。
$\theta_{\mathfrak{z}_x}$ 是满射（见 2.6.2(ii)），且与滤链严格相容
（见 2.7.4 中 $\mathfrak{z}(X)$ 上滤链的定义）。$h_X$

\SourcePage{35}

是单射，且与滤链严格相容（见 2.7.4 末尾）。
于是 $h_x$ 与滤链严格相容（这正是定理 2.5.5），
并且 $h_X$ 诱导滤代数 $\mathfrak{z}(X)$ 与
$\mathfrak{z}^{(x)}(X):=h_x(\mathfrak{z}_x)$ 之间的同构。

\bdsec{2.8. 水平性} 本小节我们在若干自然的模概形上引入 $\mathcal{D}_X$-结构，
并（模去 4.4.14 中解释的某些细节）证明水平性定理 2.7.3。
读者眼下可以跳过本小节。

在 2.8.1--2.8.2 中我们给出定理 2.7.3 的证明概要。2.8.2 的方法在 2.8.3 中略作改动。
在 2.8.4--2.8.5 中我们解释一些细节，其余细节引向 4.4.14。
在 2.8.6 中我们极简地讨论分歧情形。

\bdp{2.8.1} 让我们构造定理 2.7.3 中的态射 $h$。

回忆 2.5.3--2.5.4 中 $h_x$ 的构造涉及概形 $\operatorname{Bun}_{G,x}$，
即在 $x$ 的形式邻域上平凡化的 $X$ 上 $G$-丛的模概形。
它还涉及 Harish-Chandra 对 $(\mathfrak{g}\otimes K_x,G(\mathcal{O}_x))$
在 $\operatorname{Bun}_{G,x}$ 上的作用，以及该作用提升为
$(\widetilde{\mathfrak{g}\otimes K_x},G(\mathcal{O}_x))$ 在线丛
$\pi_x^*\omega_{\operatorname{Bun}_G}$ 上的作用，
其中 $\pi_x$ 是自然态射 $\operatorname{Bun}_{G,x}\to\operatorname{Bun}_G$。
这些作用来自群归纳概形 $G(K_x)$ 在 $\operatorname{Bun}_{G,x}$ 上的作用，
以及该作用提升为某个中心扩张\footnote{这一扩张在 2.5.3 的注中被提到（而非被定义）。
这对我们所给的概要已经足够。} $\widehat{G(K_x)}$ 在
$\pi_x^*\omega_{\operatorname{Bun}_G}$ 上的作用。

为了构造 $h$，须把上述依赖于 $x$ 的对象组织成族。
用显然的方式定义 $X$ 上的概形 $M$，其纤维在 $x$ 上等于 $\operatorname{Bun}_{G,x}$。
定义 $X$ 上的群概形 $J(G)$ 与 $X$ 上的群归纳概形 $J^{mer}(G)$，
它们在 $x$ 上的纤维分别是 $G(\mathcal{O}_x)$ 与 $G(K_x)$。
$J(G)$ 是函数 $X\to G$ 的射流概形，$J^{mer}(G)$ 是「亚纯射流」的归纳概形。
$J^{mer}(G)$ 作用于 $M$。最后，

\SourcePage{36}

定义一个中心扩张 $\widehat{J}^{mer}(G)$ 及它在
$\pi^*\omega_{\operatorname{Bun}_G}$ 上的作用，
其中 $\pi$ 是自然态射 $M\to\operatorname{Bun}_G$。
这些数据定义好之后，$h:\mathfrak{z}\to\Gamma(\operatorname{Bun}_G,D')_X$
的构造与 $h_x$ 的构造（见 2.5.3--2.5.4）颇为相似。

\bdp{2.8.2} 关键的观察是：\emph{在 $J(G)$、$J^{mer}(G)$、$\widehat{J}^{mer}(G)$、
$M$ 与 $\pi^*\omega_{\operatorname{Bun}_G}$ 上都存在沿 $X$ 的典范联络，
使得 $J^{mer}(G)$ 在 $M$ 上的作用与 $\widehat{J}^{mer}(G)$ 在
$\pi^*\omega_{\operatorname{Bun}_G}$ 上的作用都是水平的。}
由此可得 $h$ 的水平性。

对 $X$-概形 $Y$，记 $\mathcal{J}Y$ 为截面 $X\to Y$ 的射流概形。
众所周知（2.6.3 中有粗略解释），$\mathcal{J}Y$ 带有沿 $X$ 的典范联络
（即 $\mathcal{J}Y$ 是 2.6.1 意义下的 $\mathcal{D}_X$-概形）。
特别地，这适用于 $J(G)=\mathcal{J}(G\times X)$。
若 $\mathcal{F}$ 是 $X$ 上的主 $G$-丛，则 $\pi:M\to\operatorname{Bun}_G$
在 $\mathcal{F}$ 上的纤维等于 $\mathcal{J}\mathcal{F}$，故它带有沿 $X$ 的联络。
可以证明，这些联络来自 $M$ 上沿 $X$ 的联络。

为定义 $M$ 上的联络以及其它联络，使用 Grothendieck 的进路 [Gr68] 较为方便。
按 [Gr68]，$X$-概形 $Z$ 上沿 $X$ 的联络（=可积联络=「分层」）
是一族双射 $\varphi_{\alpha\beta}:\operatorname{Mor}_\alpha(S,Z)\xrightarrow{\sim}
\operatorname{Mor}_\beta(S,Z)$，对每个概形 $S$ 与每对无限接近的「点」
$\alpha,\beta:S\to X$（这里 $\operatorname{Mor}_\alpha(S,Z)$ 是 $\alpha$ 在
$\operatorname{Mor}(S,Z)$ 中的原像，「无限接近」指 $\alpha$ 与 $\beta$ 在
$S_{red}$ 上的限制相同）；双射 $\varphi_{\alpha\beta}$ 要求关于 $S$ 具有函子性，
并满足方程 $\varphi_{\beta\gamma}\varphi_{\alpha\beta}=\varphi_{\alpha\gamma}$。

例如，若 $Z$ 是 $X$ 上概形 $Y$ 的射流概形，则
$\operatorname{Mor}_\alpha(S,Z):=\operatorname{Mor}_X(S'_\alpha,Y)$，
其中 $S'_\alpha$ 是图像 $\Gamma_\alpha\subset S\times X$ 的形式邻域，
态射 $S'_\alpha\to X$ 由投影 $\operatorname{pr}_X:S\times X\to X$ 诱导。
容易证明，若 $\alpha$ 与 $\beta$ 无限接近，则 $S'_\alpha=S'_\beta$，
于是我们得到 $Z$ 上沿 $X$ 的联络。可以证明它与 2.6.3 中定义的联络一致。

\SourcePage{37}

$J^{mer}(G)$、$\widehat{J}^{mer}(G)$ 与 $M$ 上沿 $X$ 的联络用类似方式定义。
$J^{mer}(G)$ 在 $M$ 上的作用与 $\widehat{J}^{mer}(G)$ 在
$\pi^*\omega_{\operatorname{Bun}_G}$ 上的作用的水平性，由定义即可推出。

\bdp{2.8.3} 2.8.2 中描述的方法可作如下改动。回忆 $\mathcal{O}:=\mathbb{C}[[t]]$，
$K:=\mathbb{C}((t))$；$\operatorname{Aut}\mathcal{O}$ 与 $X^\wedge$ 在 2.6.5 中定义。
置 $M^\wedge=M\times_X X^\wedge$。于是 $M^\wedge$ 是四元组
$(x,t_x,\mathcal{F},\gamma_x)$ 的模空间，其中 $x\in X$，$t_x$ 是 $x$ 处的形式参数，
$\mathcal{F}$ 是 $X$ 上的 $G$-torsor，$\gamma_x$ 是 $\mathcal{F}$ 在 $x$
的形式邻域上的截面。群归纳概形 $G(K)$ 作用在 $M^\wedge$ 于任意
$\hat{x}\in X^\wedge$ 上的纤维（事实上该纤维等于 $\operatorname{Bun}_{G,x}$，
其中 $x$ 是 $\hat{x}$ 在 $X$ 中的像，故 $G(K_x)$ 作用在纤维上；
另一方面，与 $\hat{x}$ 对应的 $x$ 处形式参数给出同构 $K_x\xrightarrow{\sim}K$）。
实际上 $G(K)$ 作用在 $M^\wedge$ 上（见 2.8.4），且中心扩张 $\widehat{G(K)}$
作用在 $\hat{\pi}^*\omega_{\operatorname{Bun}_G}$ 上，
其中 $\hat{\pi}$ 是自然态射 $M^\wedge\to\operatorname{Bun}_G$。
该作用诱导态射
\[
  \hat{h}:\mathfrak{z}_{\mathfrak{g}}(\mathcal{O})\longrightarrow
  \Gamma(X^\wedge,\mathcal{O}_{X^\wedge})\otimes\Gamma(\operatorname{Bun}_G,D')
\]
（$\mathfrak{z}_{\mathfrak{g}}(\mathcal{O})$ 的定义见 2.7.2）。

另一方面，2.6.5 中 $\operatorname{Aut}\mathcal{O}$ 在 $X^\wedge$ 上的作用
典范地提升为其在 $M^\wedge$ 上（见 2.8.4）与在层
$\hat{\pi}^*\omega_{\operatorname{Bun}_G}$ 上的作用。
$\operatorname{Aut}\mathcal{O}$ 与 $\widehat{G(K)}$ 在
$\hat{\pi}^*\omega_{\operatorname{Bun}_G}$ 上的作用以显然的方式相容。
因此 $\hat{h}$ 是 $\operatorname{Aut}\mathcal{O}$-等变的。
于是 $\hat{h}$ 诱导水平态射
$h:\mathfrak{z}=\mathfrak{z}_{\mathfrak{g}}(\mathcal{O})_X\to
\Gamma(\operatorname{Bun}_G,D')_X$。

\bdp{2.8.4} 为把 2.8.3 的概要变成定理 2.7.3 的证明，
我们首先给出半直积 $\operatorname{Aut}\mathcal{O}\ltimes G(K)$ 在 $M^\wedge$
上作用的精确定义。设 $R$ 是 $\mathbb{C}$-代数。按定义，$M^\wedge$ 的 $R$-点
是一个三元组 $(\alpha,\mathcal{F},\gamma)$，其中
$\alpha:\operatorname{Spec} R\widehat{\otimes}\mathcal{O}\to X\otimes R$
是 $R$-态射，其微分在 $\operatorname{Spec} R$ 的每一点处都不为零；
$\mathcal{F}$ 是 $X\otimes R$ 上的 $G$-torsor；$\gamma$ 是 $\alpha^*\mathcal{F}$ 的截面。
记 $\Gamma_\alpha$ 为复合
$\operatorname{Spec} R\to\operatorname{Spec} R\widehat{\otimes}\mathcal{O}
\xrightarrow{\alpha}X\otimes R$ 的图像，$\alpha'$ 为 $\alpha$ 诱导的态射
$\operatorname{Spec} R\widehat{\otimes}K\to(X\otimes R)\setminus\Gamma_\alpha$。按

\SourcePage{38}

Beauville 与 Laszlo\footnote{$\Gamma_\alpha\subset X\otimes R$ 的法丛平凡，
故按 2.3.7 可以应用 [BLa95] 的主定理，而不必用其整体化版本。}
（见 2.3.7 与 2.3.4），$M^\wedge$ 的 $R$-点与三元组 $(\alpha,\mathcal{F}',\gamma')$
一一对应，其中 $\alpha$ 如上，$\mathcal{F}'$ 是 $(X\otimes R)\setminus\Gamma_\alpha$
上的 $G$-torsor，$\gamma'$ 是 $\alpha'^*\mathcal{F}'$ 的截面
（当然，$\mathcal{F}'$ 是 $\mathcal{F}$ 的限制，$\gamma'$ 是 $\gamma$ 的限制）。
这一解释表明 $G(R\widehat{\otimes}K)$ 与
$\operatorname{Aut}(R\widehat{\otimes}\mathcal{O})$ 作用在 $M^\wedge(R)$ 上：
$G(R\widehat{\otimes}K)$ 的作用改变 $\gamma'$，
$\operatorname{Aut}(R\widehat{\otimes}\mathcal{O})$ 的作用改变 $\alpha$
（若把 $\alpha$ 换成 $\alpha\circ\varphi$，
$\varphi\in\operatorname{Aut}\operatorname{Spec} R\widehat{\otimes}\mathcal{O}$，
则 $\Gamma_\alpha$ 作为 $X\otimes R$ 的子概形改变，但作为子集不变，
故 $(X\otimes R)\setminus\Gamma_\alpha$ 不变）。
这样我们得到 2.8.3 中提到的 $\operatorname{Aut}\mathcal{O}\ltimes G(K)$
在 $M^\wedge$ 上的作用。

\bdp{2.8.5} 按 2.8.4，$\operatorname{Aut}\mathcal{O}$ 作用在作为
$\operatorname{Bun}_G$ 上概形的 $M^\wedge$ 上。于是 $\operatorname{Aut}\mathcal{O}$
作用在 $\hat{\pi}^*\omega_{\operatorname{Bun}_G}$ 上。
在 2.5.3 中我们提到 $\widehat{G(K_x)}$ 在 $\omega_{\operatorname{Bun}_G}$
到 $\operatorname{Bun}_{G,x}$ 的拉回上的典范作用。
于是 $\widehat{G(K)}$ 作用在 $\hat{\pi}^*\omega_{\operatorname{Bun}_G}$
于任意 $\hat{x}\in X^\wedge$ 上纤维的限制上。如 2.8.3 中所解释，
为完成 2.7.3 的证明，只需证明：
\begin{enumerate}[label=(\roman*),leftmargin=2.6em]
\item 对应于不同 $\hat{x}\in X^\wedge$ 的 $\widehat{G(K)}$ 作用，
  来自 $\widehat{G(K)}$ 在 $\hat{\pi}^*\omega_{\operatorname{Bun}_G}$ 上的一个
  （显然唯一的）作用；
\item 该作用与 $\operatorname{Aut}\mathcal{O}$ 的作用相容。
\end{enumerate}
为证明 (i) 与 (ii)，必须（也几乎充分）定义中心扩张 $\widehat{G(K_x)}$
及其在 $\omega_{\operatorname{Bun}_G}$ 到 $\operatorname{Bun}_{G,x}$
的拉回上的作用。感兴趣的读者可以用例如 [BLa94] 来做。

我们不证明 (i) 与 (ii)，而是在 4.4.14 中对
$\omega_{\operatorname{Bun}_G}$ 的平方根证明类似的陈述
（因为我们需要 $\omega_{\operatorname{Bun}_G}$ 的平方根来表述并证明定理 5.4.5，
那是本工作的主要结果）。更精确地说，对 $\omega_X$ 的任意平方根
$\mathcal{L}$，定义 $\operatorname{Bun}_G$ 上的线丛
$\lambda'_{\mathcal{L}}$，它本质上是 $\omega_{\operatorname{Bun}_G}$ 的平方根
（见 4.4.1）。构造中心扩张\footnote{事实上，该扩张是
$\widehat{G(K_x)}$ 的平方根。} $\widehat{G(K_x)}_{\mathcal{L}}$，
它作用在 $\lambda_{\mathcal{L}}$ 到 $\operatorname{Bun}_{G,x}$ 的拉回上
（见 4.4.7--4.4.8）。从 2.5.4 来的态射
$h_x:\mathfrak{z}_x\to\Gamma(\operatorname{Bun}_G,D')$ 可以

\SourcePage{39}

用该作用自然地定义（见 4.4.12--4.4.13）。
最后，在 4.4.14 中我们对 $\lambda'_{\mathcal{L}}$ 证明上述陈述 (i) 与 (ii)
的类似物，由此推出水平性定理 2.7.3。

\bdp{2.8.6} 设 $\Delta\subset X$ 是有限子概形。记
$\operatorname{Bun}_{G,\Delta}$ 为 $X$ 上在 $\Delta$ 上平凡化的 $G$-丛的栈。
记 $D'$ 为 1.1.6 中的层 $D_{\mathcal{Y},\mathcal{L}^{\lambda}}$，
其中 $\mathcal{Y}=\operatorname{Bun}_{G,\Delta}$，
$\mathcal{L}=\omega_{\operatorname{Bun}_G}$ 的拉回，$\lambda=1/2$。
正如 $\Delta=\emptyset$ 的情形，定义水平态射
\[
  h:\mathfrak{z}_{X\setminus\Delta}\longrightarrow
  \Gamma(\operatorname{Bun}_{G,\Delta},D')\otimes\mathcal{O}_{X\setminus\Delta},
\]
其中 $\mathfrak{z}_{X\setminus\Delta}$ 是 $\mathfrak{z}$ 到 $X\setminus\Delta$ 的限制。
$h$ 诱导单射 $\Gamma(N,\mathcal{O}_N)\to\Gamma(\operatorname{Bun}_{G,\Delta},D')$，
其中 $N=N_\Delta(G)$ 是归纳概形 $N'_\Delta(G)$ 的闭子概形，
而 $N'_\Delta(G)$ 由 $\operatorname{Spec}\mathfrak{z}_{X\setminus\Delta}$
的水平截面构成。

\emph{问题.} 显式描述 $N_\Delta(G)$。

我们将指出该问题的解法中用到的几何对象。
由于不解释解法的细节，读者可以在不知道问题答案的情况下读本小节其余部分，
答案可在 3.8.2 中找到。

对 $n\in\mathbb{Z}_+$，记 $M_{\Delta,n}$ 为三元组的栈，
三元组由一个点 $x\in X$、$X$ 上的 $G$-丛 $\mathcal{F}$、
以及 $\mathcal{F}$ 在 $\Delta+nx$ 上的平凡化组成
（这里我们把 $X$ 的有限子概形与有效除子等同，故 $\Delta+nx$ 有意义）。
$M_{\Delta,n}$ 是代数栈，且
$M_\Delta:=\varprojlim_n M_{\Delta,n}$ 是 $X$ 上的概形。

\emph{注.} 设 $M_{\Delta,x}$ 是 $M_\Delta$ 在 $x\in X$ 上的纤维。
若 $x\in X\setminus\Delta$，则 $M_{\Delta,x}$ 是在 $\Delta$ 与 $x$ 的形式邻域上
平凡化的 $G$-丛的模概形。若 $x\in\Delta$，则
$M_{\Delta,x}=M_{\Delta\setminus\{x\},x}$。

考虑如下定义的「同余子群」概形 $\underline{\mathrm{G}}_\Delta$。
它是平坦于 $X$ 的概形，且对任意平坦于 $X$ 的概形 $S$ 有
\[
  \operatorname{Mor}_X(S,\underline{\mathrm{G}}_\Delta)
  =\{\,f:S\to G \text{ 满足 } f|_{\Delta_S}=1\,\},
\]
其中 $\Delta_S$ 是 $\Delta$ 在 $S$ 中的原像。
$\underline{\mathrm{G}}_\Delta$ 是 $X$ 上的群概形。
$X$ 上在 $\Delta$ 上平凡化的 $G$-丛等同于
$\underline{\mathrm{G}}_\Delta$-丛（这

\SourcePage{40}

在把 $G$-丛与 $\underline{\mathrm{G}}_\Delta$-丛看作平展拓扑下的 torsor 时
就清楚了）。于是 $\operatorname{Bun}_{G,\Delta}$ 是
$\underline{\mathrm{G}}_\Delta$-丛的栈。

可以证明，若 $D\subset X$ 是有限子概形，且 $\Delta+D$ 按除子意义理解，
则对每个平坦于 $X$ 的概形 $S$
\[
  \operatorname{Mor}_X(S,\underline{\mathrm{G}}_{\Delta+D})
  =\{\,f\in\operatorname{Mor}_X(S,\underline{\mathrm{G}}_\Delta)
  \text{ 满足 } f|_{D_S}=1\,\}.
\]

因此，$X$ 上在 $\Delta+D$ 上平凡化的 $G$-丛等同于在 $D$ 上平凡化的
$\underline{\mathrm{G}}_\Delta$-丛。于是 $M_\Delta$ 是三元组的模概形，
三元组由一个点 $x\in X$、$X$ 上的 $\underline{\mathrm{G}}_\Delta$-丛、
以及它在 $x$ 的形式邻域上的平凡化组成。
现在容易定义 $\mathcal{J}^{mer}(\underline{\mathrm{G}}_\Delta)$
在 $M_\Delta$ 上的典范作用，其中
$\mathcal{J}^{mer}(\underline{\mathrm{G}}_\Delta)$ 是截面
$X\to\underline{\mathrm{G}}_\Delta$ 的「亚纯射流」的群归纳概形。
$\mathcal{J}^{mer}(\underline{\mathrm{G}}_\Delta)$ 与 $M_\Delta$ 都带有
沿 $X$ 的联络，且上述作用是水平的。如此等等……

\emph{注}

\begin{enumerate}[label=(\roman*),leftmargin=2.6em]
\item 若 $\Delta\neq\emptyset$，则 2.8.3 的方法不能避免使用 $X$ 上的群归纳概形。
\item 与无穷维概形及归纳概形（如 $M_\Delta$ 或
  $\mathcal{J}^{mer}(\underline{\mathrm{G}}_\Delta)$）相关的陷阱是存在的。
  这里有一例。态射
  $\underline{\mathrm{G}}_\Delta\to\underline{\mathrm{G}}
  :=\underline{\mathrm{G}}_\emptyset=G\times X$ 诱导
  $f:\mathcal{J}^{mer}(\underline{\mathrm{G}}_\Delta)\to
  \mathcal{J}^{mer}(\underline{\mathrm{G}})$。
  该 $f$ 诱导任意点 $x\in X$ 上纤维的同构
  （$\mathcal{J}^{mer}(\underline{\mathrm{G}}_\Delta)$ 在 $x$ 上的纤维是
  $G(K_x)$，它不依赖 $\Delta$）。但若 $\Delta\neq\emptyset$，
  则 $f$ 不是同构，甚至不是单态射。
\end{enumerate}

\bdsec{2.9. \texorpdfstring{$\mathfrak{z}_{\mathfrak{g}}(\mathcal{O})$}{z(O)} 的交换性}
代数 $\mathfrak{z}_{\mathfrak{g}}(\mathcal{O})$ 与
$\mathfrak{z}_x=\mathfrak{z}_{\mathfrak{g}}(\mathcal{O}_x)$ 在 2.5.1 与 2.7.2
中定义（它们当然是同构的）。Feigin 与 Frenkel 在 [FF92] 中证明了
$\mathfrak{z}_{\mathfrak{g}}(\mathcal{O})$ 交换。
本小节我们给出 $\mathfrak{z}_{\mathfrak{g}}(\mathcal{O})$ 交换性的两个证明：
整体的一个（见 2.9.1--2.9.2）与局部的一个（见 2.9.3--2.9.5）。
后者实际上是 [FF92] 中原证明的一个版本。

\SourcePage{41}

读者眼下可以跳过本小节。本文其余部分不会用到 2.9.1--2.9.2。

\bdp{2.9.1} 让我们证明
\begin{equation}\tag{29}
  [h_x(\mathfrak{z}_x),h_y(\mathfrak{z}_y)]=0
\end{equation}
（$h_x:\mathfrak{z}_x\to\Gamma(\operatorname{Bun}_G,D')$ 的定义见 2.5.4）。
由于 $\mathfrak{z}_x$ 是 $\mathcal{O}_X$-代数
$\mathfrak{z}=\mathfrak{z}_{\mathfrak{g}}(\mathcal{O})_X$ 在 $x$ 处的纤维，
而 $h_x$ 来自 $\mathcal{O}_X$-代数态射
$h:\mathfrak{z}\to\mathcal{O}_X\otimes\Gamma(\operatorname{Bun}_G,D')$，
只需对 $x\neq y$ 证明 (29)。记 $\operatorname{Bun}_{G,x,y}$ 为 $X$ 上在 $x$ 与 $y$
的形式邻域上平凡化的 $G$-丛的模概形。$G(K_x)\times G(K_y)$ 作用在
$\operatorname{Bun}_{G,x,y}$ 上。特别地，Harish-Chandra 对
$\big((\mathfrak{g}\otimes K_x)\times(\mathfrak{g}\otimes K_y),
G(\mathcal{O}_x)\times G(\mathcal{O}_y)\big)$ 作用在
$\operatorname{Bun}_{G,x,y}$ 上。该作用典范地提升为
$\big(\widetilde{\mathfrak{g}\otimes K_x}\times\widetilde{\mathfrak{g}\otimes K_y},
G(\mathcal{O}_x)\times G(\mathcal{O}_y)\big)$
在 $\omega_{\operatorname{Bun}_G}$ 到 $\operatorname{Bun}_{G,x,y}$
的拉回上的作用，其中 $\mathbf{1}_x\in\mathfrak{g}\otimes K_x$ 与
$\mathbf{1}_y\in\mathfrak{g}\otimes K_y$ 按乘以 2 作用，
$G(\mathcal{O}_x)\times G(\mathcal{O}_y)$ 按显然方式作用。
$G(\mathcal{O}_x)\times G(\mathcal{O}_y)$ 在 $\operatorname{Bun}_{G,x,y}$
上的作用满足 1.2.6 中的条件 (16)，且商栈等于 $\operatorname{Bun}_G$。
于是按 1.2.5 有典范态射
$h_{x,y}:\mathfrak{z}_x\otimes\mathfrak{z}_y\to
\Gamma(\operatorname{Bun}_G,D')$。
它在 $\mathfrak{z}_x$ 与 $\mathfrak{z}_y$ 上的限制分别等于 $h_x$ 与 $h_y$。
于是 (29) 显然成立。

\bdp{2.9.2} 让我们证明 $\mathfrak{z}_{\mathfrak{g}}(\mathcal{O})$ 的交换性。
设 $a\in[\mathfrak{z}_{\mathfrak{g}}(\mathcal{O}),\mathfrak{z}_{\mathfrak{g}}(\mathcal{O})]$，
$a\neq 0$。若 $x=y$，则 (29) 表明 $h_x(\mathfrak{z}_x)$ 交换。
于是对任意 $X$、$x\in X$ 与同构 $f:\mathcal{O}\xrightarrow{\sim}\mathcal{O}_x$，
有 $h_x(f_*(a))=0$。设
$\bar{a}\in\mathfrak{z}^{cl}_{\mathfrak{g}}(\mathcal{O})$ 是 $a$ 的主象征。
则对任意如上 $X$、$x$、$f$，有 $h_x^{cl}(f_*(\bar{a}))=0$
（$h_x^{cl}:\mathfrak{z}^{cl}_x\to\Gamma(\operatorname{Bun}_G,P)
=\Gamma(T^*\operatorname{Bun}_G,\mathcal{O})$ 的定义与几何描述见 2.4.3）。
这意味着，把 $\bar{a}$ 看作 $\mathfrak{g}^*\otimes\omega_{\mathcal{O}}$
上的多项式函数（见 2.4.1）时，它具有如下性质：对任意如上 $X$、$x$，
任意在 $x$ 的形式邻域上平凡化的 $X$ 上 $G$-丛 $\mathcal{F}$，
以及任意同构 $\mathcal{O}_x\xrightarrow{\sim}\mathcal{O}$，
$\bar{a}$ 在映射
\[
  H^0(X,\mathfrak{g}^*_{\mathcal{F}}\otimes\omega_X)\longrightarrow
  \mathfrak{g}^*\otimes\omega_{\mathcal{O}_x}\xrightarrow{\sim}
  \mathfrak{g}^*\otimes\omega_{\mathcal{O}}
\]
的像上的限制为零。存在 $n$，使得 $\bar{a}$ 来自
$\mathfrak{g}^*\otimes(\omega_{\mathcal{O}}/\mathfrak{m}^n\omega_{\mathcal{O}})$
上的函数，其中 $\mathfrak{m}$ 是 $\mathcal{O}$ 的极大理想。
选取 $X$ 与 $x$，使得映射

\SourcePage{42}

$H^0(X,\omega_X)\to\omega_{\mathcal{O}_x}/\mathfrak{m}^n_x\omega_{\mathcal{O}_x}$
是满射，并令 $\mathcal{F}$ 为平凡丛。则映射
$H^0(X,\mathfrak{g}^*_{\mathcal{F}}\otimes\omega_X)\to
\mathfrak{g}^*\otimes(\omega_{\mathcal{O}_x}/\mathfrak{m}^n_x\omega_{\mathcal{O}_x})$
是满射，从而 $\bar{a}=0$，即矛盾。

\emph{注.} 设 $\operatorname{Bun}_{G,\Delta}$ 是 $X$ 上在有限子概形
$\Delta\subset X$ 上平凡化的 $G$-丛的栈。为从 (29) 推出
$\mathfrak{z}_{\mathfrak{g}}(\mathcal{O})$ 的交换性，可以使用从 $\mathfrak{z}_x$
（$x\notin\Delta$）到 $\operatorname{Bun}_{G,\Delta}$ 上扭微分算子环的自然同态。
于是，不必像上面证明中那样选取 $(X,x)$，可以固定 $(X,x)$ 并取 $\Delta$ 足够大。

\bdp{2.9.3} 记 $\bar{\mathfrak{z}}$ 为完备扭泛包络代数
$\overline{U}'(\mathfrak{g}\otimes K)$ 的中心，
$K:=\mathbb{C}((t))\supset\mathbb{C}[[t]]=\mathcal{O}$。在 [FF92] 中，
Feigin 与 Frenkel 从自然同态
$f:\bar{\mathfrak{z}}\to\mathfrak{z}_{\mathfrak{g}}(\mathcal{O})$
的满射性推出 $\mathfrak{z}_{\mathfrak{g}}(\mathcal{O})$ 的交换性。
我们将给出 $f$ 满射性的一个证明：它可视为 [FF92] 中证明的几何版本，
也可视为 2.4.2 末尾注的「量子化」。
与 [FF92] 及 2.4.2 的关系将在 2.9.7 与 2.9.8 中说明。

\emph{注.} 在中心扩张 $\mathfrak{g}\otimes K$ 的定义中（见 2.5.1），
从而在 $\bar{\mathfrak{z}}$ 与 $\mathfrak{z}_{\mathfrak{g}}(\mathcal{O})$ 的定义中，
我们使用了由 (18) 定义的「临界」双线性形式 $c$。
在 $f$ 满射性的证明中，可以假设 $c$ 是 $\mathfrak{g}$ 上任一不变对称双线性形式，
且 $\mathfrak{g}$ 是任一有限维李代数。另一方面，已知若 $\mathfrak{g}$ 单且 $c$ 非临界，
则相应的代数 $\mathfrak{z}_{\mathfrak{g}}(\mathcal{O})$ 平凡（见 ???）。
% [Tutor Note: 原书 p.42 此处交叉引用印作「???」，底本如此，保留原样。]

\bdp{2.9.4} 我们需要 [BD94] 中
$\overline{U}':=\overline{U}'(\mathfrak{g}\otimes K)$ 的解释。
记 $U'$ 为 $\mathfrak{g}\otimes K$ 的非完备扭泛包络代数。
对 $n\ge 0$，设 $I_n$ 是 $U'$ 中由
$\mathfrak{g}\otimes\mathfrak{m}^n\subset\mathfrak{g}\otimes\mathcal{O}\subset U'$
生成的左理想。按定义，$\overline{U}':=\varprojlim_n U'/I_n$。
设 $U'_k$ 是 $U'$ 的标准滤链，$\overline{U}'_k$ 是 $U'_k$ 在
$\overline{U}'$ 中的闭包，即
$\overline{U}'_k:=\varprojlim_n U'_k/I_{n,k}$，$I_{n,k}:=I_n\cap U'_k$。
[BD94] 的主定理把对偶空间 $(U'_k/I_{n,k})^*$ 与某个

\SourcePage{43}

拓扑向量空间 $\Omega_{n,k}$ 等同。于是
\begin{equation}\tag{30}
  U'_k/I_{n,k}=(\Omega_{n,k})^*,\qquad \overline{U}'_k=(\Omega_k)^*
\end{equation}
其中 $\Omega_k=\varinjlim_n\Omega_{n,k}$，$*$ 表示拓扑对偶。

为定义 $\Omega_{n,k}$，需要一些记号。记 $\mathcal{O}_r$
（相应地，$\omega^{\mathcal{O}}_r$）为 $r$ 个 $\mathcal{O}$
（相应地，$\omega_{\mathcal{O}}$）的完备张量积。置
$\omega^K_r=\omega^{\mathcal{O}}_r\widehat{\otimes}_{\mathcal{O}_r}K_r$，
其中 $K_r$ 是 $\mathcal{O}_r$ 的分式域。我们把 $\mathcal{O}_r$ 与
$\mathbb{C}[[t_1,\ldots,t_r]]$ 等同，并把 $\omega^K_r$ 的元素写成
$f(t_1,\ldots,t_r)\,dt_1\ldots dt_r$，其中 $f$ 属于
$\mathbb{C}[[t_1,\ldots,t_r]]$ 的分式域。

\emph{定义.} $\Omega_{n,k}$ 是 $(k+1)$-元组 $(w_0,\ldots,w_k)$ 的集合，
其中 $w_r\in(\mathfrak{g}^*)^{\otimes r}\otimes\omega^K_r$，且满足
\begin{enumerate}[label=\arabic*),leftmargin=2.4em]
\item $w_r$ 关于对称群 $S_r$ 的作用不变（$S_r$ 同时作用在
  $(\mathfrak{g}^*)^{\otimes r}$ 与 $\omega^K_r$ 上）；
\item $w_r$ 在超平面 $t_i=0$（$1\le i\le r$）上有阶 $\le n$ 的极点，
  在超平面 $t_i=t_j$（$1\le i<j\le r$）上有阶 $\le 2$ 的极点，
  且没有其它极点；
\item 若 $w_r=f_r(t_1,\ldots,t_r)\,dt_1\ldots dt_r$，$r\ge 2$，则
\begin{equation}\tag{31}
\begin{aligned}
  f_r(t_1\ldots,t_r)&=\frac{f_{r-2}(t_1,\ldots,t_{r-2})\otimes c}
    {(t_{r-1}-t_r)^2}\\
  &\quad+\frac{\varphi^*(f_{r-1}(t_1,\ldots,t_{r-1}))}{t_{r-1}-t_r}+\cdots
\end{aligned}
\end{equation}
\end{enumerate}
这里 $c\in\mathfrak{g}^*\otimes\mathfrak{g}^*$ 是中心扩张
$\mathfrak{g}\otimes K$ 的定义中用到的双线性形式，
$\varphi^*:(\mathfrak{g}^*)^{\otimes(r-1)}\to(\mathfrak{g}^*)^{\otimes r}$
是映射 $\varphi:\mathfrak{g}^{\otimes r}\to\mathfrak{g}^{\otimes(r-1)}$ 的对偶，
后者由
$\varphi(a_1\otimes\ldots\otimes a_r)
=a_1\otimes\ldots\otimes a_{r-2}\otimes[a_{r-1},a_r]$ 给出；
而 (31) 中的点号表示在超平面 $t_{r-1}=t_r$ 的一般点处没有极点的表达式。

$\Omega_{n,k}$ 上的拓扑由嵌入
\[
  \Omega_{n,k}\hookrightarrow\prod_{0\le r\le k}
  (\mathfrak{g}^*)^{\otimes r}\otimes\Omega^{\mathcal{O}}_r
\]
诱导，该嵌入由 $(w_0,\ldots,w_k)\mapsto(\eta_0,\ldots,\eta_k)$ 给出，其中
\[
  \eta_r=\prod_i t_i^n\cdot\prod_{i<j}(t_i-t_j)^2\cdot w_r.
\]

% ch02_u04: 原书 p.44--p.52
\SourcePage{44}

说明一下：(31) 中我们把 $f_r$ 看作取值于 $(\mathfrak{g}^*)^{\otimes r}$ 的函数。

我们不需要 [BD94] 中同构 (30) 的显式公式。只提一点：按 [BD94] 的命题 5，
$\mathfrak{g}\otimes K$ 在 $\overline{U}'_k$ 上的伴随作用经 (30) 诱导出
$\mathfrak{g}\otimes K$ 在 $\Omega_k$ 上的如下作用：$a\in\mathfrak{g}\otimes K$
把 $(w_0,\ldots,w_k)\in\Omega_k$ 送到 $(0,w''_1,\ldots,w''_k)$，其中
\[
\begin{aligned}
w''_r&=\frac{1}{(r-1)!}\operatorname{Sym} w'_r,\\
w'_r(t_1,\ldots,t_r)&:=(\operatorname{id}\otimes\ldots\otimes\operatorname{id}
  \otimes\operatorname{ad}_{a(t_r)})w_r(t_1,\ldots,t_r)\\
&\quad-w_{r-1}(t_1,\ldots,t_{r-1})\otimes c\cdot da(t_r).
\end{aligned}
\]
这里 $\operatorname{Sym}$ 表示对称化算子（不含因子 $1/r!$），
$\operatorname{ad}_{a(t_r)}:\mathfrak{g}^*\to\mathfrak{g}^*$ 是余伴随表示中
对应于 $a(t_r)$ 的算子，$c:\mathfrak{g}\to\mathfrak{g}^*$ 是 $\mathfrak{g}$
的双线性形式。

\emph{注.} 设 $c=0$，且 $\mathfrak{g}$ 交换。则
$U'_k/I_{n,k}=\bigoplus_{r=0}^k\operatorname{Sym}^r(\mathfrak{g}\otimes K/\mathfrak{m}^n)$，
$\Omega_{n,k}=\bigoplus_{r=0}^k\overline{\operatorname{Sym}}^r
(\mathfrak{g}^*\otimes\mathfrak{m}^{-n}\omega_{\mathcal{O}})$，
其中 $\overline{\operatorname{Sym}}^r$ 表示完备对称幂。同构
$U'_k/I_{n,k}\xrightarrow{\sim}(\Omega_{n,k})^*$ 就是把
$\operatorname{Sym}(\mathfrak{g}\otimes K/\mathfrak{m}^n)$ 与
$\mathfrak{g}^*\otimes\mathfrak{m}^{-n}\omega_{\mathcal{O}}$ 上的多项式函数空间
等同起来；这个等同在 2.4.1 与 2.4.2 中用过。

\bdp{2.9.5} 按 2.9.4，要证
$f:\bar{\mathfrak{z}}\to\mathfrak{z}_{\mathfrak{g}}(\mathcal{O})$ 满射，
只需证明：任一 $(\mathfrak{g}\otimes\mathcal{O})$-不变连续线性泛函
$l:\Omega_{0,k}\to\mathbb{C}$ 都可延拓为 $(\mathfrak{g}\otimes K)$-不变连续线性泛函
$\Omega_k\to\mathbb{C}$。考虑连续线性算子
\[
T:\Omega_k\to\mathbb{C}((\zeta))\widehat{\otimes}\Omega_{0,k}
=\Bigl\{\sum_{n=-\infty}^{\infty}a_n\zeta^n
\Bigm| a_n\in\Omega_{0,k},\ a_n\to 0\ (n\to-\infty)\Bigr\},
\]
它由下式定义：
\begin{equation}\tag{32}
T(w_0,\ldots,w_k)=(\widehat{w}_0,\ldots,\widehat{w}_k),\qquad
\widehat{w}_r=w_r(\zeta+t_1,\ldots,\zeta+t_r),
\end{equation}
其中 $w_r(\zeta+t_1,\ldots,\zeta+t_r)$ 看作
\[
\Delta^{-1}\mathbb{C}((\zeta))[[t_1,\ldots,t_n]]\,dt_1\ldots dt_n
=\mathbb{C}((\zeta))\widehat{\otimes}\Delta^{-1}\mathbb{C}[[t_1,\ldots,t_n]]
\,dt_1\ldots dt_n
\]
的元素。
% [Tutor Note: 原书 p.44 把 t_r 写成 t_n，此处照录。]

\SourcePage{45}

\[
\Delta:=\prod_{1\le i<j\le r}(t_i-t_j)^2.
\]
设 $l\in(\Omega_{0,k})^*$。令 $\bar{l}:\Omega_k\to\mathbb{C}((\zeta))$ 为
$T:\Omega_k\to\mathbb{C}((\zeta))\widehat{\otimes}\Omega_{0,k}$ 与
$\operatorname{id}\otimes l:\mathbb{C}((\zeta))\widehat{\otimes}\Omega_{0,k}
\to\mathbb{C}((\zeta))$ 的合成。把 $\bar{l}$ 写成 $\sum_i l_i\zeta^i$，
$l_i\in(\Omega_k)^*$。若 $l$ 是 $\mathfrak{g}\otimes\mathcal{O}$-不变的，
则各泛函 $l_i$ 是 $\mathfrak{g}\otimes K$-不变的。此外
$l_0|_{\Omega_{0,k}}=l$。

\emph{注.} 设 $G$ 是代数群，且 $\operatorname{Lie}G=\mathfrak{g}$。
则 $G(K)$ 作用在 $\mathfrak{g}\otimes K$ 的中心扩张上（见 (19)），
从而作用在 $\overline{U}'_k$ 上；此外 $G(\mathcal{O})$ 作用在 $U'_k/I_{n,k}$ 上。
因此 $G(K)$ 作用在 $\Omega_k$ 上，$G(\mathcal{O})$ 作用在 $\Omega_{n,k}$ 上。
在上述情形下，若 $l$ 是 $G(\mathcal{O})$-不变的，则各 $l_i$ 是 $G(K)$-不变的
（$G(K)$ 在 $\Omega_k$ 上的作用见 [BD94] 的公式 (24)）。注意：即使 $G$ 连通，
$G(K)$ 也未必连通。因此 $G(K)$-不变性不能由 $(\mathfrak{g}\otimes K)$-不变性
直接推出。

\bdp{2.9.6} 由于 $\bar{l}$ 连续，当 $i\to-\infty$ 时 $l_i\to 0$
（即对每个 $n$，只要 $i$ 足够大就有 $l_{-i}(\Omega_{n,k})=0$）。
于是映射 $l\mapsto\bar{l}$ 可看作从 $U'_k/I_{0,k}$ 到
$W_k:=\bigl\{\sum_{i=-\infty}^{\infty}a_i\zeta^i \bigm| a_i\in\overline{U}'_k,\ 
a_i\to 0\ (i\to-\infty)\bigr\}$ 的映射。这些映射定义一个算子
\begin{equation}\tag{33}
\Phi:\mathit{Vac}'\to W:=\bigcup_k W_k
\end{equation}
其中 $\mathit{Vac}'=U'/I_0$ 是扭真空模。按 2.9.5 的解释，$\Phi$ 诱导映射
\begin{equation}\tag{34}
\mathfrak{z}_{\mathfrak{g}}(\mathcal{O})\to
\mathfrak{z}\widehat{\otimes}\mathbb{C}((\zeta))
:=\Bigl\{\sum_{i=-\infty}^{\infty}a_i\zeta^i
\Bigm| a_i\in\mathfrak{z},\ a_i\to 0\ (i\to-\infty)\Bigr\}.
\end{equation}
可以证明 (34) 是环同态（见 ???）。容易看出：(34) 与投影
$\mathfrak{z}\widehat{\otimes}\mathbb{C}((\zeta))
\to\mathfrak{z}_{\mathfrak{g}}(\mathcal{O})((\zeta))$ 的合成把
$\mathfrak{z}_{\mathfrak{g}}(\mathcal{O})$ 映入
$\mathfrak{z}_{\mathfrak{g}}(\mathcal{O})[[\zeta]]$；并且合成
$\mathfrak{z}_{\mathfrak{g}}(\mathcal{O})\to\mathfrak{z}_{\mathfrak{g}}(\mathcal{O})[[\zeta]]
\xrightarrow{\zeta=0}\mathfrak{z}_{\mathfrak{g}}(\mathcal{O})$ 是恒等映射。
% [Tutor Note: 原书 p.45 此处交叉引用印作「???」，底本如此，保留原样。]

\emph{注.} 设 $G$ 是连通代数群，且 $\operatorname{Lie}G=\mathfrak{g}$。
则 (34) 的像中所有元素都是 $G(K)$-不变的（见 2.9.5 的注）。

\SourcePage{46}

\bdp{2.9.7} 可以证明，(33) 与 Feigin 和 Frenkel 构造的算子
$F:\mathit{Vac}'\to W$ 一致（见 [FF92] 引理 1 的证明）。因此 2.9.5
只是 [FF92] 中一部分内容的一个版本。

[FF92] 中 $F$ 的定义可改写如下。置
$W_k^+:=\overline{U}'_k((\zeta))$；$W_k^-$ 由所有满足「当 $i$ 足够大时
$a_{-i}=0$」的 $\sum_i a_i\zeta^i\in W_k$ 组成。
定义 $W^\pm\subset W$ 为 $W^\pm=\bigcup_k W_k^\pm$。$W^+$ 与 $W^-$
有自然的代数结构，$W$ 有自然的 $(W^+,W^-)$-双模结构
（$W$ 是左 $W^+$-模，也是右 $W^-$-模）。考虑线性映射
$\varphi^\pm:\mathfrak{g}\widehat{\otimes}K\to W^\pm$，使得
\[
\varphi^+(1)=1,\qquad \varphi^-(1)=0,
\]
且对 $a\in\mathfrak{g}((t))=\mathfrak{g}\otimes K
\subset\mathfrak{g}\widehat{\otimes}K$，
\[
\varphi^+(a)=a(t-\zeta)\in\mathfrak{g}((t))((\zeta)),\qquad
\varphi^-(a)=a(t-\zeta)\in\mathfrak{g}((\zeta))((t)).
\]
容易证明 $\varphi^\pm$ 是李代数同态。考虑 $W$ 上的
$\mathfrak{g}\widehat{\otimes}K$-模结构
$a\circ w:=\varphi_+(a)w-w\varphi_-(a)$，
$a\in\mathfrak{g}\widehat{\otimes}K$，$w\in W$。则 $F:\mathit{Vac}'\to W$
是 $\mathfrak{g}\widehat{\otimes}K$-模同态，它把 $\mathit{Vac}'$ 的真空向量
送到 $1\in W$。

\bdp{2.9.8} 说明 (34) 与 2.4.2 中其经典类比的关系。

$\overline{U}'$ 带有标准滤链 $\overline{U}'_k$（见 2.9.4）。它诱导滤链
$\mathfrak{z}_k:=\mathfrak{z}\cap\overline{U}_k$。
% [Tutor Note: 原书 p.46 此处印作 \overline{U}_k（无撇），按 2.9.4 的记号应为 \overline{U}'_k，此处照录。]
我们把 $\operatorname{gr}_k\overline{U}':=
\overline{U}'_k/\overline{U}'_{k-1}$ 与
$\operatorname{Sym}^k(\mathfrak{g}\otimes K)$ 的完备化等同，即与 $k$
次齐次多项式函数 $\mathfrak{g}^*\otimes\omega_K\to\mathbb{C}$ 的空间等同，
其中 $\omega_K:=\omega_{\mathcal{O}}\otimes_{\mathcal{O}}K$（称
$\mathfrak{g}^*\otimes\omega_K$ 上的函数 $f$ 为多项式函数，是指：对每个
$n$，它在 $\mathfrak{g}^*\otimes\mathfrak{m}^{-n}$ 上的限制是多项式，即来自
$\mathfrak{g}^*\otimes(\mathfrak{m}^{-n}/\mathfrak{m}^N)$ 上的多项式函数，
$N$ 依赖于 $n$）。记 $\mathfrak{z}^{cl}$ 为 $\mathfrak{g}^*\otimes\omega_K$ 上
$\mathfrak{g}\otimes K$-不变多项式函数的代数。显然
$\operatorname{gr}\mathfrak{z}$ 在 $\operatorname{gr}\overline{U}'$ 中的像
含于 $\mathfrak{z}^{cl}$。

\SourcePage{47}

$\mathfrak{z}$ 的滤链诱导 $\mathfrak{z}\widehat{\otimes}\mathbb{C}((\zeta))$
上的滤链，且映射 (34) 与滤链相容。我们断言下图交换：
\begin{equation}\tag{35}
\begin{tikzcd}[column sep=4.6em, row sep=2.8em]
\operatorname{gr}_k\mathfrak{z}_{\mathfrak{g}}(\mathcal{O}) \arrow[r] \arrow[d, "\sigma"]
  & \operatorname{gr}_k\mathfrak{z}\widehat{\otimes}\mathbb{C}((\zeta)) \arrow[d] \\
\mathfrak{z}^{cl}_{\mathfrak{g}}(\mathcal{O}) \arrow[r, "\nu"]
  & \mathfrak{z}^{cl}\widehat{\otimes}\mathbb{C}((\zeta))
\end{tikzcd}
\end{equation}
这里上方的箭头由 (34) 诱导，$\mathfrak{z}^{cl}_{\mathfrak{g}}(\mathcal{O})$
的定义见 2.4.1，$\sigma$ 是 1.2.5 中的象征映射，$\nu$ 由下式定义：
\begin{equation}\tag{36}
\begin{gathered}
\nu(f):=h(\zeta),\qquad (h(\zeta))(\varphi):=f(\varphi(\zeta+t))\\
f\in\mathfrak{z}^{cl}_{\mathfrak{g}}(\mathcal{O}),\quad
\varphi\in\mathfrak{g}^*\otimes\omega_K,\quad
\varphi(\zeta+t)\in\mathfrak{g}^*((\zeta))[[t]]dt
=(\mathfrak{g}^*\otimes\omega_{\mathcal{O}})\widehat{\otimes}\mathbb{C}((\zeta)).
\end{gathered}
\end{equation}
这里把 $\mathfrak{z}^{cl}_{\mathfrak{g}}(\mathcal{O})$ 与
$\mathfrak{g}^*\otimes\omega_{\mathcal{O}}$ 上 $\mathfrak{g}\otimes\mathcal{O}$-不变
多项式函数的代数等同（参看 2.4.1）。映射 $\nu$ 在 2.4.2 的注中讨论过。

(35) 的交换性由下图 (37) 的交换性推出：
\begin{equation}\tag{37}
\begin{tikzcd}[column sep=4.6em, row sep=2.8em]
(U'_k/I_{n,k})^{*} \arrow[r, "\sim"] \arrow[d, "\sigma^{*}"]
  & \Omega_{n,k} \arrow[d] \\
(\operatorname{Sym}^k(\mathfrak{g}\otimes K/\mathfrak{g}\otimes\mathfrak{m}^n))^{*}
  \arrow[r, "\sim"]
  & ((\mathfrak{m}^{-n}\omega_{\mathcal{O}})^{\otimes k})^{S_k}
\end{tikzcd}
\end{equation}
这里上方的箭头是 (30) 的对偶，
$\sigma:U'_k/I_{n,k}\to
\operatorname{Sym}^k(\mathfrak{g}\otimes K/\mathfrak{g}\otimes\mathfrak{m}^n)$
是象征映射，右侧竖箭头由 $w\mapsto(0,\ldots,0,w)$ 定义。(37) 的交换性
由 (30) 的定义直接推出；见 [BD94]。

\bdsec{2.10. $T^{*}\operatorname{Bun}_G$ 的几何} 本小节应看作附录；
读者完全可以跳过。

置 $\operatorname{Nilp}=\operatorname{Nilp}(G):=p^{-1}(0)$，其中
$p:T^{*}\operatorname{Bun}_G\to\operatorname{Hitch}(X)$ 是 Hitchin 纤维化
（见 2.2.3）。$\operatorname{Nilp}$ 由 [La87] 与 [La88] 引入，名为整体幂零锥
（若 $\mathcal{F}$ 是 $X$ 上 $G$-丛，且
$\eta\in T^{*}_{\mathcal{F}}\operatorname{Bun}_G
=H^0(X,\mathfrak{g}^{*}_{\mathcal{F}}\otimes\omega_X)$，则
$(\mathcal{F},\eta)\in\operatorname{Nilp}$ 当且仅当 $\eta$ 在
$H^0(X,\mathfrak{g}_{\mathcal{F}}\otimes\omega_X)$ 中的像是幂零的）。

\SourcePage{48}

在 2.10.1 中我们将证明：命题 2.2.4(iii) 容易由等式
\begin{equation}\tag{38}
\dim\operatorname{Nilp}=\dim\operatorname{Bun}_G
\end{equation}
推出。

我们还由 (38) 推出：$\operatorname{Bun}_G$ 在 1.1.1 的意义下是良的。
等式 (38) 由 Faltings 与 Ginzburg 证明；在 $G=PSL_n$ 的特殊情形，
它已由 Laumon 证明。在 2.10.2 中我们对其证明作一些评述。在 2.10.3 中
讨论 $\operatorname{Nilp}$ 的不可约分支集。在 2.10.4 中证明：即使 $X$
的亏格等于 $0$ 或 $1$，$\operatorname{Nilp}$ 也是等维的（若 $g>1$，
这由 2.2.4(iii) 推出）。在 2.10.5 中证明 $\operatorname{Bun}_G$ 在 1.1.1
的意义下是极良的。

我们将用 $\mathfrak{g}$ 上不变的数量积把 $\mathfrak{g}$ 与
$\mathfrak{g}^*$ 等同。

\bdp{2.10.1} 假定 (38) 成立，我们来证明 2.2.4(iii)，并证明
$\operatorname{Bun}_G$ 在 1.1.1 的意义下是良的。设
$U\subset T^{*}\operatorname{Bun}_G$ 是满足
$\dim U\le 2\dim\operatorname{Bun}_G$ 的最大开子栈。(38) 的意思是：
$p:T^{*}\operatorname{Bun}_G\to\operatorname{Hitch}(X)$ 在 $0$ 上的纤维
维数为 $\dim\operatorname{Bun}_G$。由于
$\dim\operatorname{Hitch}(X)=\dim\operatorname{Bun}_G$，这说明
$U\supset p^{-1}(0)$。$U$ 关于 $\mathbb{G}_m$ 在
$T^{*}\operatorname{Bun}_G$ 上的自然作用不变。因此
$U=T^{*}\operatorname{Bun}_G$。于是
$\dim T^{*}\operatorname{Bun}_G\le 2\dim\operatorname{Bun}_G$。
按 1.1.1，这意味着 $\operatorname{Bun}_G$ 是良的，且
$T^{*}\operatorname{Bun}_G$ 是纯维数 $2\dim\operatorname{Bun}_G$
的局部完全交。

对开子集 $V\subset T^{*}\operatorname{Bun}_G$，以下两条性质等价：
1) $p$ 在 $V$ 上的限制平坦；2) 该限制的纤维维数为
$\dim\operatorname{Bun}_G$。设 $V_{\max}$ 是具有这两条性质的最大 $V$。
$V_{\max}$ 是 $\mathbb{G}_m$-不变的，且按 (38) 有
$V_{\max}\supset p^{-1}(0)$。于是 $V_{\max}=T^{*}\operatorname{Bun}_G$，
这就证明了 2.2.4(iii) 的第一个断言。它蕴含 $p^{\gamma}$ 的像是开的。
另一方面，该像 $\mathbb{G}_m$-不变且包含 $0$。因此 $p^{\gamma}$ 是满射。
证毕。

\SourcePage{49}

由于 $\operatorname{Nilp}$ 包含 $T^{*}\operatorname{Bun}_G$ 的零截面，
(38) 可由不等式
$\dim\operatorname{Nilp}\le\dim\operatorname{Bun}_G$ 推出；后者在
[La88]、[Fal93]、[Gi97] 中作为下述定理的推论得到。

\bdp{2.10.2} \emph{定理.}（[La88]，[Fal93]，[Gi97]）
$\operatorname{Nilp}$ 是迷向的。

\emph{注.}
\begin{enumerate}[label=(\roman*),leftmargin=2.6em]
\item 说明一下：光滑辛簇 $M$ 的子概形 $N$ 称为迷向的，是指 $N$ 的任一
  光滑子簇都是迷向的。可以证明：$N$ 迷向当且仅当 $N_{\mathrm{red}}$
  的非奇异点集迷向。称 $N$ 为 Lagrange 的，是指 $N$ 迷向，且对任意
  $x\in N$ 有 $\dim_x N=1/2\dim_x M$。若 $\mathcal{Y}$ 是光滑代数栈，
  则称子栈 $\mathcal{N}\subset T^{*}\mathcal{Y}$ 为迷向的（相应地，
  Lagrange 的），是指对某个呈现\footnote{$\mathcal{Y}$ 的一个呈现是指
  光滑满态射 $S\to\mathcal{Y}$，其中 $S$ 是概形。}
  $S\to\mathcal{Y}$，有
  $\mathcal{N}\times_{\mathcal{Y}}S\subset
  (T^{*}\mathcal{Y})\times_{\mathcal{Y}}S\subset T^{*}S$
  迷向（相应地，Lagrange）；此时对一切呈现 $S\to\mathcal{Y}$ 也都成立。
\item 定理 2.10.2 在 [Fal93] 与 [Gi97] 中给出的证明都不使用 $g>1$
  的假设，这里 $g$ 是 $X$ 的亏格。若 $g>1$，Faltings 与 Ginzburg
  证明 $\operatorname{Nilp}$ 是 Lagrange 的。其论证已在 2.10.1 中解释：
  (38) 蕴含 $\operatorname{Nilp}$ 的维数纯为
  $\dim\operatorname{Bun}_G$。在 2.10.1 中我们用到了等式
  $\dim\operatorname{Hitch}(X)=\dim\operatorname{Bun}_G$，它只在 $g>1$
  时成立。事实上，即使 $g=0,1$，$\operatorname{Nilp}$ 也是 Lagrange 的
  （见 2.10.4）。
\item 由于 $\operatorname{Nilp}\subset T^{*}\operatorname{Bun}_G$ 是
  Lagrange 的且 $\mathbb{G}_m$-不变，它是一些余法丛的并，这些余法丛对应
  $\operatorname{Bun}_G$ 的某些既约不可约闭子栈。对 $G=PSL_n$，
  Laumon 给出了其中某些子栈的描述（见 [La88] 的 §§3.8--3.9）。
\item Ginzburg 对定理 2.10.2 的证明基于如下解释：用
  $\pi:\operatorname{Bun}_B\to\operatorname{Bun}_G$ 刻画
  $\operatorname{Nilp}$，其中 $B$ 是 $G$ 的 Borel 子群；即若
  $\mathcal{F}\in\operatorname{Bun}_G$，
  $\eta\in T^{*}_{\mathcal{F}}\operatorname{Bun}_G$，则
  $(\mathcal{F},\eta)\in\operatorname{Nilp}$
\SourcePage{50}
  当且仅当存在 $\mathcal{E}\in\pi^{-1}(\mathcal{F})$，使得 $\eta$ 在
  $T^{*}_{\mathcal{E}}\operatorname{Bun}_B$ 中的像等于 $0$。
  借助这一解释，Ginzburg 用辛几何中一个简单而一般的论证证明了
  定理 2.10.2（见 [Gi97] 的 §§6.5）。Faltings 对定理 2.10.2 的证明
  也很漂亮且简短（见 [Fal93] 定理 II.5 证明的前两段）。
\item {}[La88] 中关于 $G=PSL_n$ 的定理 2.10.2 证明在一般情形不成立，
  因为它用到 $\mathfrak{g}=sl_n$ 的如下性质：对每个幂零元
  $A\in\mathfrak{g}$，存在抛物子群 $P\subset G$，使得 $A$ 属于幂幺根
  $U\subset P$ 的李代数，$A$ 的 $P$-轨道在 $\operatorname{Lie}U$ 中
  开，且 $A$ 在 $G$ 中的中心化子含于 $P$。该性质对 $\mathfrak{g}=sl_n$
  成立（例如可取 $P$ 为旗 $0\subset\operatorname{Ker}A
  \subset\operatorname{Ker}A^2\subset\cdots$ 的稳定化子），但对任意
  半单 $\mathfrak{g}$ 不成立（例如当 $\mathfrak{g}=sp_4$、$A\in sp_4$
  是秩 $1$ 的幂零算子时不成立）。
\end{enumerate}

\bdp{2.10.3} 本小节我们「描述」$\operatorname{Nilp}$ 的不可约分支集。

回忆：$\operatorname{Nilp}$ 是二元组 $(\mathcal{F},\eta)$ 构成的栈，
其中 $\mathcal{F}$ 是 $X$ 上 $G$-丛，
$\eta\in H^0(X,\mathfrak{g}_{\mathcal{F}}\otimes\omega_X)
=H^0(X,\mathfrak{g}^{*}_{\mathcal{F}}\otimes\omega_X)$ 幂零。对幂零共轭类
$C\subset\mathfrak{g}$，有局部闭子栈 $\operatorname{Nilp}_C$，
它参数化满足「对一般 $x\in X$ 有 $\eta(x)\in C$」的二元组
$(\mathcal{F},\eta)$。

固定 $e\in C$，把它扩充为 $sl_2$-三元组 $\{e,f,h\}$。设
$\mathfrak{g}^k$ 是 $\mathfrak{g}$ 的降滤链，使得
$[h,\mathfrak{g}^k]\subset\mathfrak{g}^k$，且 $\operatorname{ad}_h$
在 $\mathfrak{g}^k/\mathfrak{g}^{k+1}$ 上以乘以 $k$ 的方式作用。
$\mathfrak{g}^k$ 依赖 $e$，但不依赖 $h$ 与 $f$。置
$\mathfrak{p}=\mathfrak{p}_e:=\mathfrak{g}^0$。$\mathfrak{p}$ 是
$\mathfrak{g}$ 的抛物子代数。设 $P\subset G$ 是相应的子群。有映射
$C\to G/P$，它把 $a\in C$ 对应到抛物子代数 $\mathfrak{p}_a$。其纤维
$\{a\in C\mid\mathfrak{p}_a=\mathfrak{p}\}$（即 $e\in C$ 的 $P$-轨道）
等于 $\mathfrak{g}^2\cap C$；这是 $\mathfrak{g}^2$ 的开子集。称
$\mathfrak{g}^2$ 的元素为一般的，是指它属于 $\mathfrak{g}^2\cap C$。

\SourcePage{51}

设 $(\mathcal{F},\eta)\in\operatorname{Nilp}_C$，
$U:=\{x\in X\mid\eta(x)\in C\}$。$\eta\in
\Gamma(U,C_{\mathcal{F}}\otimes\omega_X)$ 在
$\Gamma(U,(G/P)_{\mathcal{F}})$ 中的像可延拓为 $(G/P)_{\mathcal{F}}$
在 $X$ 上的截面。于是得到 $\mathcal{F}$ 上的 $P$-结构。就这个
$P$-结构而言，$\eta\in H^0(X,\mathfrak{g}^2_{\mathcal{F}}\otimes\omega_X)$，
且对 $x\in U$，$\eta(x)$ 是一般的。

记 $Y_C$ 为二元组 $(\mathcal{F},\eta)$ 构成的栈，其中 $\mathcal{F}$
是 $X$ 上 $P$-丛，$\eta\in
H^0(X,\mathfrak{g}^2_{\mathcal{F}}\otimes\omega_X)$，且对几乎所有
$x\in X$，$\eta(x)$ 是一般的。对 $P$-丛 $\mathcal{F}$，记
$\deg\mathcal{F}\in\operatorname{Hom}(P,\mathbb{G}_m)^*$ 为如下泛函：
它把 $\varphi:P\to\mathbb{G}_m$ 对应到 $\mathcal{F}$ 经 $\varphi$
的前推的次数。$Y_C$ 是开子栈 $Y_C^u$
（$u\in\operatorname{Hom}(P,\mathbb{G}_m)^*$）的无交并，$Y_C^u$
参数化满足 $\deg\mathcal{F}=u$ 的 $(\mathcal{F},\eta)\in Y_C$。
容易证明：对每个 $u\in\operatorname{Hom}(P,\mathbb{G}_m)^*$，
自然态射 $Y_C^u\to\operatorname{Nilp}_C$ 是局部闭嵌入，且子栈
$Y_C^u\subset\operatorname{Nilp}_C$ 构成 $\operatorname{Nilp}_C$
的一个分层。

\emph{引理.}
\begin{enumerate}[label=\arabic*),leftmargin=2.2em]
\item $Y_C^u$ 是光滑等维栈。
  $\dim Y_C^u\le\dim\operatorname{Bun}_G$。
\item 设 $Y_C^*$ 是 $Y_C$ 的维数为 $\dim\operatorname{Bun}_G$ 的连通分支
  之并。则 $Y_C^*$ 是满足下述条件的二元组 $(\mathcal{F},\eta)\in Y_C$
  构成的栈：$\operatorname{ad}_\eta:
  (\mathfrak{g}^{-1}/\mathfrak{g}^0)_{\mathcal{F}}\to
  (\mathfrak{g}^1/\mathfrak{g}^2)_{\mathcal{F}}\otimes\omega_X$
  是同构。
\end{enumerate}

\emph{注.} (38) 由引理推出。

\emph{证明.} 二元组 $(\mathcal{F},\eta)\in Y_C^u$ 的形变理论由复形
$C^{\bullet}$ 的超上同调控制，其中
$C^0=\mathfrak{p}_{\mathcal{F}}=\mathfrak{g}^0_{\mathcal{F}}$，
$C^1=\mathfrak{g}^2_{\mathcal{F}}\otimes\omega_X$，$i\ne 0,1$ 时
$C^i=0$，微分 $d:C^0\to C^1$ 等于 $\operatorname{ad}_\eta$。由于
$\operatorname{Coker}d$ 有有限支撑，
$\mathbb{H}^2(X,C^{\bullet})=0$。因此 $Y_C$ 光滑，且
\[
\begin{aligned}
\dim_{(\mathcal{F},\eta)}Y_C
&=\chi(\mathfrak{g}^2_{\mathcal{F}}\otimes\omega_X)
  -\chi(\mathfrak{g}^0_{\mathcal{F}})
 =-\chi(\mathfrak{g}_{\mathcal{F}}/\mathfrak{g}^{-1}_{\mathcal{F}})
  -\chi(\mathfrak{g}^0_{\mathcal{F}})\\
&=-\chi(\mathfrak{g}_{\mathcal{F}})
  +\chi(\mathfrak{g}^{-1}_{\mathcal{F}}/\mathfrak{g}^0_{\mathcal{F}})
 =\dim\operatorname{Bun}_G
  +\chi(\mathfrak{g}^{-1}_{\mathcal{F}}/\mathfrak{g}^0_{\mathcal{F}}).
\end{aligned}
\]
显然 $\chi(\mathfrak{g}^{-1}_{\mathcal{F}}/\mathfrak{g}^0_{\mathcal{F}})$
只依赖 $u=\deg\mathcal{F}$。态射 $\operatorname{ad}_\eta:
\mathfrak{g}^{-1}_{\mathcal{F}}/\mathfrak{g}^0_{\mathcal{F}}\to
(\mathfrak{g}^1/\mathfrak{g}^2)_{\mathcal{F}}\otimes\omega_X$ 是单射，
其余核 $\mathcal{A}$ 有有限支撑。于是
$2\chi(\mathfrak{g}^{-1}_{\mathcal{F}}/\mathfrak{g}^0_{\mathcal{F}})
=\chi(\mathfrak{g}^{-1}_{\mathcal{F}}/\mathfrak{g}^0_{\mathcal{F}})
-\chi((\mathfrak{g}^1/\mathfrak{g}^2)_{\mathcal{F}}\otimes\omega_X)
=-\chi(\mathcal{A})\le 0$，且
$\chi(\mathfrak{g}^{-1}_{\mathcal{F}}/\mathfrak{g}^0_{\mathcal{F}})=0$
当且仅当 $\mathcal{A}=0$。\hfill$\Box$

\SourcePage{52}

由于 $\operatorname{Nilp}$ 的维数纯为 $\dim\operatorname{Bun}_G$，
引理蕴含：$\operatorname{Nilp}$ 的不可约分支由
$\bigsqcup_C\pi_0(Y_C^*)$ 参数化。

$\pi_0(Y_C^*)$ 可等同于一个更简单的栈 $M_C$ 的 $\pi_0$，后者定义如下。
置 $L=P/U$，其中 $U$ 是 $P$ 的幂幺根。$L$ 作用在
$V:=\mathfrak{g}^2/\mathfrak{g}^3$ 上。记 $D_i$ 为满足下述条件的
$a\in V$ 的集合：$(\operatorname{ad}_a)^i:
\mathfrak{g}^{-i}/\mathfrak{g}^{-i+1}\to
\mathfrak{g}^i/\mathfrak{g}^{i+1}$ 的行列式为 $0$。
$D_i\subset V$ 是纯余维 $1$ 的 $L$-不变闭子集。$\mathfrak{g}^2$
的元素一般，当且仅当它在 $V$ 中的像不属于 $D_2$。因此对所有 $i$ 有
$D_i\subset D_2$。记 $M_C$ 为二元组 $(\mathcal{F},\eta)$ 构成的栈，
其中 $\mathcal{F}$ 是 $X$ 上 $L$-丛，
$\eta\in H^0(X,V_{\mathcal{F}}\otimes\omega_X)$ 满足：对一切
$x\in X$ 有 $\eta(x)\notin D_1$，且对一般的 $x\in X$ 有
$\eta(x)\notin D_2$。容易看出，自然态射 $Y_C^*\to M_C$ 诱导双射
$\pi_0(Y_C^*)\to\pi_0(M_C)$。

于是 $\operatorname{Nilp}$ 的不可约分支由
$\bigsqcup_C\pi_0(M_C)$ 参数化。但愿 $\pi_0(M_C)$ 能用与 $C$ 和
$X$ 相伴的「标准」对象来描述……

\emph{注.} 若 $G=PSL_n$，则对每个幂零共轭类 $C\subset sl_n$，
$\operatorname{Nilp}_C$ 的维数纯为 $\dim\operatorname{Bun}_G$
（见 [La88]）。这并非总成立。例如取 $G=Sp_4$，$C$ 为 $sp_4$ 中秩
$1$ 的幂零矩阵之集。事实上，设 $(\mathcal{F},\eta)\in Y_C$ 使得
$\eta\in H^0(X,\mathfrak{g}^2_{\mathcal{F}}\otimes\omega_X)$ 只有单零点。
则容易证明：态射 $Y_C\to\operatorname{Nilp}_C$ 在
$(\mathcal{F},\eta)$ 的邻域内是开嵌入。另一方面，由上述引理，若
$\eta$ 有零点，则 $Y_C$ 在 $(\mathcal{F},\eta)$ 处的维数小于
$\operatorname{Bun}_G$。

\bdp{2.10.4} \emph{定理.} $\operatorname{Nilp}$ 是 Lagrange 的。

本定理不假设 $g>1$。

\emph{证明.} 按 2.10.2 的注 (ii)，只需证明当 $g\le 1$ 时
$\operatorname{Nilp}$ 的维数纯为 $\dim\operatorname{Bun}_G$。

1) 设 $g=0$。则 $\operatorname{Nilp}=T^{*}\operatorname{Bun}_G$。
$\operatorname{Bun}_G$ 的拟紧开子栈可表为 $H\backslash M$，其中 $M$
是光滑簇，$H$ 是
% [Tutor Note: 原书 p.52 在此处断页，句子未完，下接 p.53。]
%% begin units/ch02_u05.tex
\SourcePage{53}代数群，作用在 $M$ 上。则 $T^{*}(H\backslash M)=H\backslash N$，其中
$N\subset T^{*}M$ 是 $H$ 各轨道之余法丛的并。每个余法丛的维数纯为
$\dim M$；又因 $g=0$，$H$-轨道的个数有限。

\emph{注.} 本质上相同的论证表明：对任何光滑代数栈 $\mathcal{Y}$，
$T^{*}\mathcal{Y}$ 在每一点的维数都 $\ge\dim\mathcal{Y}$。若
$g=0$ 且 $\mathcal{Y}=\operatorname{Bun}_G$，则按定理 2.10.2 有
$T^{*}\mathcal{Y}=\operatorname{Nilp}$，且
$\dim T^{*}\mathcal{Y}=\dim\mathcal{Y}$。于是又证得 $g=0$ 时的
定理 2.10.4。

2) 设 $g=1$。此时宜设 $G$ 约化，而不必半单（这其实并非本质：
约化 $G$ 的定理 2.10.4 易由半单情形推出）。

在进入证明之前，先回忆半稳定性与 Shatz 分层的概念。固定 Borel 子群
$B\subset G$，记 $H$ 为其极大交换商。设 $P\subset G$ 是含 $B$ 的
抛物子群，$L$ 为 $P$ 的极大约化商，$Z$ 为 $L$ 的中心。设 $\Gamma$
（resp. $\Delta$）为 $G$（resp. $L$）的单根集。嵌入
$Z\hookrightarrow L$ 诱导同构
\[
  \operatorname{Hom}(Z,\mathbb{G}_m)\otimes\mathbb{Q}
  \xrightarrow{\sim}
  \operatorname{Hom}(L,\mathbb{G}_m)\otimes\mathbb{Q}.
\]
记 $p$ 为合成
\[
  \operatorname{Hom}(H,\mathbb{G}_m)\to
  \operatorname{Hom}(Z,\mathbb{G}_m)\to
  \operatorname{Hom}(L,\mathbb{G}_m)\otimes\mathbb{Q}
  =\operatorname{Hom}(P,\mathbb{G}_m)\otimes\mathbb{Q}.
\]
称 $l\in\operatorname{Hom}(P,\mathbb{G}_m)^{*}$ \emph{严格支配}，是指对
$\alpha\in\Gamma\backslash\Delta$ 有 $l(p(\alpha))>0$。

对 $P$-丛 $\mathcal{F}$，记
$\deg\mathcal{F}\in\operatorname{Hom}(P,\mathbb{G}_m)^{*}$ 为如下泛函：
它把 $\varphi:P\to\mathbb{G}_m$ 对应到 $\mathcal{F}$ 经 $\varphi$
的前推的次数。称 $G$-丛为\emph{半稳定}的，是指对任何 $P\ne G$，
它都不来自严格支配次数的 $P$-丛。半稳定的 $G$-丛构成开子栈
$\operatorname{Bun}_G^{ss}\subset\operatorname{Bun}_G$。固定次数
$d\in\operatorname{Hom}(G,\mathbb{G}_m)$ 的半稳定 $G$-丛构成开子栈
$\operatorname{Bun}_G^{ss,d}\subset\operatorname{Bun}_G^{ss}$。若
$P\subset G$ 是含 $B$ 的抛物子群，
$d\in\operatorname{Hom}(P,\mathbb{G}_m)^{*}$ 严格支配，则记
$\operatorname{Shatz}_P^d$ 为次数为 $d$ 且相应 $L$-丛半稳定的
$P$-丛 $\mathcal{F}$ 之栈。已知自然态射
$\operatorname{Shatz}_P^d\to\operatorname{Bun}_G$ 是局部闭嵌入，且子栈%
\SourcePage{54}%
$\operatorname{Shatz}_P^d$（$P,d$ 遍取）构成 $\operatorname{Bun}_G$
的一个分层，称为 \emph{Shatz 分层}。

记 $\operatorname{Nilp}_P^d(G)$（resp. $\operatorname{Nilp}^{ss}(G)$、
$\operatorname{Nilp}^{ss,d}(G)$）为
$\operatorname{Nilp}=\operatorname{Nilp}(G)$ 与
$\operatorname{Shatz}_P^d$（resp. $\operatorname{Bun}_G^{ss}$、
$\operatorname{Bun}_G^{ss,d}$）在 $\operatorname{Bun}_G$ 上的纤维积。
要证 $\operatorname{Nilp}(G)$ 的维数纯为
$\dim\operatorname{Bun}_G=0$，只需对每个 $P,d$ 证明
$\operatorname{Nilp}_P^d(G)$ 的维数纯为 $0$。设 $L$ 为 $P$ 的极大
约化商，$\mathfrak{p}:=\operatorname{Lie}P$，
$\mathfrak{l}:=\operatorname{Lie}L$。若 $\mathcal{F}$ 是具有严格支配
次数的 $P$-丛，且相应的 $L$-丛 $\overline{\mathcal{F}}$ 半稳定，则
$H^0(X,\mathfrak{g}_{\mathcal{F}})
=H^0(X,\mathfrak{p}_{\mathcal{F}})$，于是有自然映射
$\eta\mapsto\overline{\eta}$，从
$H^0(X,\mathfrak{g}_{\mathcal{F}})$ 到
$H^0(X,\mathfrak{l}_{\mathcal{F}})$。定义
$\pi:\operatorname{Nilp}_P^d(G)\to\operatorname{Nilp}^{ss,d}(L)$ 为
$(\mathcal{F},\eta)\mapsto(\overline{\mathcal{F}},\overline{\eta})$，
其中 $\eta\in H^0(X,\mathfrak{g}_{\mathcal{F}}\otimes\omega_X)
=H^0(X,\mathfrak{g}_{\mathcal{F}})$（因 $g=1$，$\omega_X$ 平凡）。
再次利用 $g=1$ 可知 $\pi$ 光滑，其纤维是 $0$ 维栈。故只需证
$\operatorname{Nilp}^{ss}(L)$ 的维数纯为 $0$。

$\operatorname{Nilp}^{ss}(L)$ 的点是二元组：半稳定的 $L$-丛
$\mathcal{F}$ 与幂零元 $\eta\in H^0(X,\mathfrak{l}_{\mathcal{F}})$。
因 $\mathfrak{l}_{\mathcal{F}}$ 是半稳定向量丛，
$\operatorname{ad}_\eta:\mathfrak{l}_{\mathcal{F}}\to
\mathfrak{l}_{\mathcal{F}}$ 的秩恒定。故 $\eta(x)$ 的共轭类不依赖
$x\in X$。对幂零共轭类 $C\subset\mathfrak{l}$，记
$\operatorname{Nilp}_C^{ss}(L)$ 为 $\operatorname{Nilp}^{ss}(L)$ 中
参数化满足 $\eta(x)\in C$ 的二元组 $(\mathcal{F},\eta)$ 的局部闭
子栈。只需对每个 $C$ 证明 $\operatorname{Nilp}_C^{ss}(L)$ 的维数
纯为 $0$。设 $Z(A)\subset L$ 是某 $A\in C$ 的中心化子，
$\mathfrak{z}(A):=\operatorname{Lie}Z(A)$。若
$(\mathcal{F},\eta)\in\operatorname{Nilp}_C^{ss}(L)$，则
$\eta\in\Gamma(X,C_{\mathcal{F}})
=\Gamma(X,(G/Z(A))_{\mathcal{F}})$ 在 $\mathcal{F}$ 上定义
$Z(A)$-结构。于是得到开嵌入
$\operatorname{Nilp}_C^{ss}(L)
\hookrightarrow\operatorname{Bun}_{Z(A)}$。最后，
$\operatorname{Bun}_{Z(A)}$ 的维数纯为 $0$：对任何 $Z(A)$-丛
$\mathcal{E}$ 有 $\chi(\mathfrak{z}(A)_{\mathcal{E}})
=\deg\mathfrak{z}(A)_{\mathcal{E}}=0$（注意 $G/Z(A)=C$ 带 $G$-不变
辛结构，故 $Z(A)$ 的伴随表示行列式平凡，从而
$\mathfrak{z}(A)_{\mathcal{E}}$ 平凡）。\hfill$\Box$

\bdp{2.10.5} \emph{命题 2.1.2 的证明.} 须证 (4) 对
$\mathcal{Y}=\operatorname{Bun}_G$ 成立，即对一切 $n>0$ 有
$\operatorname{codim}\{\mathcal{F}\in\operatorname{Bun}_G\mid
\dim H^0(X,\mathfrak{g}_{\mathcal{F}})=n\}>n$。

\SourcePage{55}

这等价于证明
\begin{equation}\tag{39}
  \dim(A(G)\backslash A^{0}(G))<\dim\operatorname{Bun}_G,
\end{equation}
其中 $A(G)$ 是二元组 $(\mathcal{F},s)$ 之栈，$\mathcal{F}\in
\operatorname{Bun}_G$，$s\in H^0(X,\mathfrak{g}_{\mathcal{F}})$；而
$A^{0}(G)\subset A(G)$ 是由方程 $s=0$ 定义的闭子栈。

置 $C:=\operatorname{Spec}(\operatorname{Sym}\mathfrak{g}^{*})^{G}$。
这是 $\mathfrak{g}$ 关于 $G$ 的伴随作用的仿射概形商；事实上
$C=W\backslash\mathfrak{h}$，其中 $\mathfrak{h}$ 是 $\mathfrak{g}$ 的
固定嘉当子代数，$W$ 是 Weyl 群。态射 $\mathfrak{g}\to C$ 诱导映射
$H^0(X,\mathfrak{g}_{\mathcal{F}})\to\operatorname{Mor}(X,C)=C$。
于是有典范态射 $f:A(G)\to C=W\backslash\mathfrak{h}$。对
$h\in\mathfrak{h}$，置 $A_h(G)=f^{-1}(\overline{h})$，其中
$\overline{h}\in W\backslash\mathfrak{h}$ 是 $h$ 的像。置
$G^{h}:=\{g\in G\mid ghg^{-1}=h\}$，$\mathfrak{g}^{h}:=
\operatorname{Lie}G^{h}=\{a\in\mathfrak{g}\mid[a,h]=0\}$。记
$\mathfrak{z}_h$ 为 $\mathfrak{g}^{h}$ 的中心。由于
$h\in\mathfrak{z}_h$，且 $\mathfrak{g}$ 中形如 $\mathfrak{z}_h$ 的
子代数只有有限个，(39) 由不等式
$\dim(A_h(G)\backslash A^{0}(G))
<\dim\operatorname{Bun}_G-\dim\mathfrak{z}_h$ 推出。故只需证明
\begin{equation}\tag{40}
  \dim A_h(G)<\dim\operatorname{Bun}_G-\dim\mathfrak{z}_h
  \qquad\text{当 } h\ne 0
\end{equation}
\begin{equation}\tag{41}
  \dim(A_0(G)\backslash A^{0}(G))<\dim\operatorname{Bun}_G.
\end{equation}

记 $Z_h$ 为 $G^{h}$ 的中心。下面说明：把 (41) 中的 $G$ 换成
$G^{h}/Z_h$，(40) 即可推出。事实上，有自然同构
\[
  A_0(G^{h})\xrightarrow{\sim}A_h(G^{h})\xrightarrow{\sim}A_h(G),
\]
以及显然的态射 $\varphi:A_0(G^{h})\to A_0(G^{h}/Z_h)$。
$\varphi$ 的非空纤维同构于 $\operatorname{Bun}_{Z_h}$，故
$\dim A_h(G)\le\dim\operatorname{Bun}_{Z_h}
+\dim A_0(G^{h}/Z_h)$。因
$\dim\operatorname{Bun}_{Z_h}=(g-1)\cdot\dim\mathfrak{z}_h$，而
(41) 蕴含 $\dim A_0(G^{h}/Z_h)
=(g-1)\cdot\dim(\mathfrak{g}^{h}/\mathfrak{z}_h)$，故
\[
  \dim A_h(G)\le(g-1)\cdot\dim\mathfrak{g}^{h}
  =\dim\operatorname{Bun}_G
  -(g-1)\cdot\dim(\mathfrak{g}/\mathfrak{g}^{h})
  \le\dim\operatorname{Bun}_G-\dim(\mathfrak{g}/\mathfrak{g}^{h}).
\]
最后，若 $h\ne 0$，则 $\dim(\mathfrak{g}/\mathfrak{g}^{h})
\ge 2\cdot\dim\mathfrak{z}_h>\dim\mathfrak{z}_h$。

为证 (41)，我们将证明：若 $Y\subset A_0(G)$ 是局部闭既约不可约
子栈，则 $\dim Y\le\dim\operatorname{Bun}_G$，且仅当
$Y\subset A^{0}(G)$ 时 $\dim Y=\dim\operatorname{Bun}_G$。对
$\xi\in H^0(X,\omega_X)$，考虑态射
$m_\xi:A_0(G)\to\operatorname{Nilp}$，
$(\mathcal{F},s)\mapsto(\mathcal{F},s\xi)$，其中 $\mathcal{F}\in
\operatorname{Bun}_G$，$s\in H^0(X,\mathfrak{g}_{\mathcal{F}})$。
这些态射%
\SourcePage{56}%
$m_\xi$ 定义 $m:A_0(G)\times H^0(X,\omega_X)\to\operatorname{Nilp}$。
$m$ 的像含于某个局部闭既约不可约子栈
$Z\subset\operatorname{Nilp}$ 中。若 $\xi\ne 0$，则 $m_\xi$ 诱导
嵌入 $Y\hookrightarrow Z_\xi$，其中 $Z_\xi$ 是 $Z$ 的闭子栈，由满足
下述条件的二元组 $(\mathcal{F},\eta)\in
H^0(X,\mathfrak{g}_{\mathcal{F}}\otimes\omega_X)$ 组成：$\eta$ 在
子概形 $D_\xi:=\{x\in X\mid\xi(x)=0\}$ 上的限制为零。故
$\dim Y\le\dim Z_\xi\le\dim Z\le\dim\operatorname{Nilp}
=\dim\operatorname{Bun}_G$。若
$\dim Y=\dim\operatorname{Bun}_G$，则对所有非零
$\xi\in H^0(X,\omega_X)$ 有 $Z_\xi=Z$。这意味着所有
$(\mathcal{F},\eta)\in Z$ 都有 $\eta=0$，从而所有
$(\mathcal{F},s)\in Y$ 都有 $s=0$，即 $Y\subset A^{0}(G)$。
\hfill$\Box$

\bdsec{2.11. 论局部系统之栈} 记 $\mathcal{LS}_G$ 为 $X$ 上
$G$-局部系统之栈（$G$-局部系统即带联络的 $G$-丛）。Kapranov
[Kap97] 说明 $\mathcal{LS}_G$ 有导出版本 $R\mathcal{LS}_G$，后者是
DG 栈。利用 2.10 的结果，我们将证明：若 $g>1$ 且 $G$ 半单，则
$R\mathcal{LS}_G=\mathcal{LS}_G$。我们还将描述 $\mathcal{LS}_G$ 的
不可约分支之集。读者可跳过本节；其余部分不采用本节的结果。

\bdp{2.11.1} 固定 $x\in X$。记 $\mathcal{LS}_G^x$ 为 $X$ 上带联络
$\nabla$ 的 $G$-丛 $\mathcal{F}$ 之栈，该联络在 $x$ 有单极点。
% [Tutor Note: 原书 p.56 此处印作「G-biundles」，疑为「G-bundles」之笔误，底本如此，保留原样。]
记 $\mathcal{E}$ 为万有 $G$-丛（它定义在
$\mathcal{LS}_G^x\times X$ 上）在
$\mathcal{LS}_G^x=\mathcal{LS}_G^x\times\{x\}$ 上的限制。
% [Tutor Note: 原书 p.56 此处等式左端多写 $\mathcal{LS}_G^x$（应为 $\mathcal{LS}_G^x\times\{x\}$），底本如此，保留原样。]
$\nabla$ 在 $x$ 处的留数是截面
$R\in\Gamma(\mathcal{LS}_G^x,\mathfrak{g}_{\mathcal{E}})$，而
$\mathcal{LS}_G$ 是 $\mathcal{LS}_G^x$ 中由方程 $R=0$ 定义的闭子栈。
考虑开子栈 $\widetilde{\mathcal{LS}}{}_G^x\subset\mathcal{LS}_G^x$，
它参数化使 $\nabla:H^1(X,\mathfrak{g}_{\mathcal{F}})\to
H^1(X,\mathfrak{g}_{\mathcal{F}}\otimes\omega_X(x))$ 满射的二元组
$(\mathcal{F},\nabla)$。易见 $\widetilde{\mathcal{LS}}{}_G^x$ 是
纯维数 $(2g-1)\cdot\dim G$ 的光滑栈，且
$\mathcal{LS}_G\subset\widetilde{\mathcal{LS}}{}_G^x$。

把 $\mathfrak{g}_{\mathcal{E}}$ 看作 $\mathcal{LS}_G^x$ 上的栈。截面
$R,0\in\Gamma(\mathcal{LS}_G^x,\mathfrak{g}_{\mathcal{E}})$ 定义
$\mathfrak{g}_{\mathcal{E}}$ 的两个闭子栈；$R\mathcal{LS}$ 是它们在
导出意义下的交，而 $\mathcal{LS}_G$ 是通常的交。故下列条件等价：
\begin{enumerate}[label=\arabic*),leftmargin=2.2em]
\item $R\mathcal{LS}_G=\mathcal{LS}_G$；
\SourcePage{57}
\item $\mathcal{LS}_G$ 是纯维数 $(2g-2)\cdot\dim G$ 的局部完全交；
\item $\dim\mathcal{LS}_G\le(2g-2)\cdot\dim G$。
\end{enumerate}

下述命题表明：若 $g>1$ 且 $G$ 半单，则这些条件都满足。

\bdp{2.11.2} \emph{命题.} 设 $g>1$，$G$ 约化。则 $\mathcal{LS}_G$
是纯维数 $(2g-2)\cdot\dim G+l$ 的局部完全交，其中 $l$ 是 $G$ 的
中心的维数。

\emph{证明.} 设 $R$ 的含义与 2.11.1 中相同。显然
$R\in\Gamma(\mathcal{LS}_G^x,[\mathfrak{g},\mathfrak{g}]_{\mathcal{E}})$，
故只需证明
\begin{equation}\tag{42}
  \dim\mathcal{LS}_G\le(2g-2)\cdot\dim G+l.
\end{equation}
记 $G_{ad}$ 为 $G$ 模去中心所得的商。考虑投影
$p:\mathcal{LS}_G\to\operatorname{Bun}_{G_{ad}}$。若 $p$ 在
$G_{ad}$-丛 $\mathcal{F}$ 上的纤维非空，则其维数等于
$\dim T^{*}_{\mathcal{F}}\operatorname{Bun}_{G_{ad}}+l(2g-1)$，故
$\dim\mathcal{LS}_G\le\dim T^{*}\operatorname{Bun}_{G_{ad}}+l(2g-1)$。
最后，$\dim T^{*}\operatorname{Bun}_{G_{ad}}
\le\dim G_{ad}\cdot(2g-2)$，因为 $\operatorname{Bun}_{G_{ad}}$ 在
1.1.1 的意义下是良的（2.10.1 中已证）。\hfill$\Box$

\bdp{2.11.3} 记 $\operatorname{Bun}'_G\subset\operatorname{Bun}_G$
为 $\operatorname{Bun}_{G/[G,G]}$ 中含平凡丛的连通分支的原像。
$\mathcal{LS}_G\to\operatorname{Bun}_G$ 的像含于
$\operatorname{Bun}'_G$。

\bdp{2.11.4} \emph{命题.} 设 $g>1$，$G$ 约化。则
$\operatorname{Bun}'_G$ 的每个连通分支在 $\mathcal{LS}_G$ 中的原像
都非空且不可约。

于是 $\mathcal{LS}_G$ 的不可约分支由
\[
  \operatorname{Ker}(\pi_1(G)\to\pi_1(G/[G,G]))=\pi_1([G,G])
\]
参数化。

\emph{证明.} 考虑开子栈
$\operatorname{Bun}^0_{G_{ad}}\subset\operatorname{Bun}_{G_{ad}}$，
它参数化满足 $H^0(X,(\mathfrak{g}_{ad})_{\mathcal{F}})=0$ 的
$G_{ad}$-丛 $\mathcal{F}$（这是 $\operatorname{Bun}_{G_{ad}}$ 中
最大的 Deligne--Mumford 子栈）。记 $\operatorname{Bun}^{0\prime}_G$
为 $\operatorname{Bun}^0_{G_{ad}}$ 在 $\operatorname{Bun}'_G$ 中的
原像。

\SourcePage{58}

记 $\mathcal{LS}^0_G$ 为 $\operatorname{Bun}^0_{G_{ad}}$ 在
$\mathcal{LS}_G$ 中的原像。2.10.5 中已证
$\operatorname{Bun}_{G_{ad}}$ 在 1.1.1 的意义下是极良的，故
$\dim(T^{*}\operatorname{Bun}_{G_{ad}}
\backslash T^{*}\operatorname{Bun}^0_{G_{ad}})
<\dim T^{*}\operatorname{Bun}_{G_{ad}}$。(42) 的证明中所用的论证
表明 $\dim(\mathcal{LS}_G\backslash\mathcal{LS}^0_G)
<(2g-2)\cdot\dim G+l$。由 2.11.2 可知 $\mathcal{LS}^0_G$ 在
$\mathcal{LS}_G$ 中稠密。故只需证明
$\operatorname{Bun}^{0\prime}_G$ 的每个连通分支在
$\mathcal{LS}^0_G$ 中的原像都非空且不可约。这显然：态射
$\mathcal{LS}^0_G\to\operatorname{Bun}^{0\prime}_G$ 是
$T^{*}\operatorname{Bun}^{0\prime}_G$ 上的
torsor\footnote{torsor 结构依赖 $\mathfrak{g}$ 上一个不变数量积的
选取。}。\hfill$\Box$

\bdsec{2.12. 论 Beauville--Laszlo 定理} 本节实际上是一个附录，其中
解释 [BLa95] 主定理的整体化版本。2.3.7 用到这个版本，但并非本质地
用到。故读者可跳过本节。

\bdp{2.12.1} \emph{定理.} 设 $p:\widetilde{S}\to S$ 是概形态射，
$D\subset S$ 是有效 Cartier 除子。设
$\widetilde{D}:=p^{-1}(D)$ 是 $\widetilde{S}$ 中的 Cartier 除子，
且态射 $\widetilde{D}\to D$ 是同构。置 $U:=S\backslash D$，
$\widetilde{U}:=\widetilde{S}\backslash\widetilde{D}$。

记 $C$ 为由这样的拟凝聚 $\mathcal{O}_S$-模构成的范畴：它没有支撑在
$D$ 上的非零局部截面。记 $\widetilde{C}$ 为
$(\widetilde{S},\widetilde{D})$ 的类似范畴。记 $C'$ 为三元组
$(\mathcal{M}_1,\mathcal{M}_2,\varphi)$ 的范畴，其中
$\mathcal{M}_1$ 是拟凝聚 $\mathcal{O}_U$-模，
$\mathcal{M}_2\in\widetilde{C}$，$\varphi$ 是 $\mathcal{M}_1$ 与
$\mathcal{M}_2$ 到 $\widetilde{U}$ 的拉回之间的同构。
\begin{enumerate}[label=\arabic*),leftmargin=2.2em]
\item $p^{*}$ 把 $C$ 送到 $\widetilde{C}$，故有函子
  $F:C\to C'$，它把 $\mathcal{M}\in C$ 送到
  $(\mathcal{M}|_U,p^{*}\mathcal{M},\varphi)$，其中 $\varphi$ 是
  $\mathcal{M}|_U$ 与 $p^{*}\mathcal{M}$ 到 $\widetilde{U}$ 的拉回
  之间的自然同构。
\item $F:C\to C'$ 是范畴等价。
\item $\mathcal{M}\in C$ 局部有限型（resp. 平坦，resp. 有限秩局部
  自由），当且仅当 $\mathcal{M}|_U$ 与 $f^{*}\mathcal{M}$ 有此性质。
% [Tutor Note: 原书 p.58 此处印作 $f^{*}\mathcal{M}$；本节只定义了 $p$，疑为 $p^{*}\mathcal{M}$ 之笔误，底本如此，保留原样。]
\end{enumerate}

\SourcePage{59}

这一定理容易化归为下述情形：$S$ 与 $\widetilde{S}$
仿射\footnote{对任意 $x\in S$，存在 $x$ 的仿射邻域 $U$ 与开仿射
$\widetilde{U}\subset\widetilde{S}$，使得
$\widetilde{U}\subset p^{-1}(U)$，且 $\widetilde{U}\cap\widetilde{D}
=p^{-1}(U)\cap\widetilde{D}$。事实上，可设 $S$ 仿射且 $x\in D$。
设 $\widetilde{U}_1\subset\widetilde{S}$ 是 $x$ 在
$\widetilde{D}$ 中的原像的仿射邻域。则
$p(\widetilde{U}_1\cap\widetilde{D})$ 是 $x$ 在 $D$ 中的仿射邻域，
故对某个满足 $x\in U$ 的开仿射 $U\subset S$，它含 $U\cap D$。于是
$\widetilde{U}:=\widetilde{U}_1\times_S U$ 具有所需性质。}，且 $D$ 由
单个方程整体地定义（即 $S=\operatorname{Spec}A$，
$\widetilde{S}=\operatorname{Spec}\widetilde{A}$，
$D=\operatorname{Spec}A/fA$，$f\in A$ 不是零因子）。这种情形与
[BLa95] 中的处理相同（[BLa95] 假设
$\widetilde{A}=\widehat{A}:=$ $A$ 关于 $f$-进拓扑的完备化，但
[BLa95] 只用到 $\widehat{A}$ 的两条性质：$f:\widehat{A}\to
\widehat{A}$ 的单射性，以及 $A/fA\to\widehat{A}/f\widehat{A}$ 的
双射性）。

\bdp{2.12.2} 设 $D$ 是概形 $S$ 的闭仿射子概形。记 $\widehat{S}$ 为
$S$ 沿 $D$ 的完备化，$\widehat{S}'$ 为 $\widehat{S}$ 上正则函数环的
谱（故 $\widehat{S}$ 是仿射形式概形，$\widehat{S}'$ 是相应的真
概形）。有态射 $\pi:\widehat{S}\to S$ 与
$i:\widehat{S}\to\widehat{S}'$。

\bdp{2.12.3} \emph{命题.} 至多存在一个态射
$p:\widehat{S}'\to S$ 使 $pi=\pi$。

\emph{证明.} 设对某 $p_1,p_2:\widehat{S}'\to S$ 有
$\pi=p_1i=p_2i$。设 $Y\subset\widehat{S}'$ 是对角
$\Delta\subset S\times S$ 在
$(p_1,p_2):\widehat{S}'\to S\times S$ 下的原像。则 $Y$ 是
$\widehat{S}'$ 的局部闭子概形，且对每个 $n$ 都含
$D\subset\widehat{S}'$ 的第 $n$ 个无穷小邻域。故
$(\overline{Y}\backslash Y)\cap D=\varnothing$，从而
$\overline{Y}\backslash Y=\varnothing$，即 $Y$ 闭。含 $D$ 的所有
无穷小邻域的 $\widehat{S}'$ 闭子概形等于 $\widehat{S}'$。故
$Y=\widehat{S}'$，$p_1=p_2$。\hfill$\Box$

\bdp{2.12.4} 设处于 2.12.2 的情形，且 $D\subset S$ 是有效 Cartier
除子。若存在 $p:\widehat{S}'\to S$ 使 $pi=\pi$，则
$p^{-1}(D)\subset\widehat{S}'$ 是 Cartier 除子，且态射
$p^{-1}(D)\to D$ 是同构。于是定理 2.12.1 适用。

\SourcePage{60}

\bdp{2.12.5} 设处于 2.12.2 的情形，且 $S$ 拟分离。则存在
$p:\widehat{S}'\to S$ 使 $pi=\pi$。我们所知的证明相当长：先处理
诺特情形，再用下述事实（Deligne 私下告知）：对任何拟紧拟分离概形
$S$，存在从 $S$ 到某个 $\mathbb{Z}$ 上有限型概形的仿射态射。

在 2.3.7 中，我们对 $S=X\otimes R$ 用到 $p:\widehat{S}'\to S$ 的
存在性，其中 $X$ 是我们的曲线，$R$ 是 $\mathbb{C}$-代数。故下述
结果够用。

\bdp{2.12.6} \emph{命题.} 设在 2.12.2 的情形中，$S$ 是
$\mathbb{P}^{n}\otimes R$ 的局部闭子概形（$R$ 为某环）。则存在
$p:\widehat{S}'\to S$ 使 $pi=\pi$。

\emph{证明.} 我们用 Jouanolou 技巧。设 $\mathbb{P}^{*}$ 是与
$\mathbb{P}=\mathbb{P}^{n}$ 对偶的射影空间，
$Z\subset\mathbb{P}\times\mathbb{P}^{*}$ 为关联对应，
$U:=(\mathbb{P}\times\mathbb{P}^{*})\backslash Z$。因态射
$U\to\mathbb{P}$ 是 $\mathbb{P}$ 上某向量丛上的 torsor，而
$\widehat{S}$ 是仿射形式概形，态射 $\widehat{S}\to\mathbb{P}$ 提升
为态射 $\widehat{S}\to U$。因 $U$ 仿射，
$\operatorname{Mor}(\widehat{S},U)=\operatorname{Mor}(\widehat{S}',U)$，
故得态射 $\widehat{S}'\to U$。合成
$\widehat{S}'\to U\to\mathbb{P}$ 给出态射
$f:\widehat{S}'\to\mathbb{P}\otimes R$。局部闭子概形
$f^{-1}(S)\subset\widehat{S}'$ 对每个 $n$ 都含
$D\subset\widehat{S}'$ 的第 $n$ 个无穷小邻域，故
$f^{-1}(S)=\widehat{S}'$（参看 2.12.3），且 $f$ 诱导态射
$p:\widehat{S}'\to S\subset\mathbb{P}\otimes R$。显然
$pi=\pi$。\hfill$\Box$

\emph{注.} 也可用下述方式证明本命题：把态射
$\widehat{S}\to\mathbb{P}^{n}$ 解释为二元组
$(\mathcal{M},\varphi)$，其中 $\mathcal{M}$ 是 $\widehat{S}$ 上的
可逆层，$\varphi$ 是满态射 $\mathcal{O}^{n+1}\to\mathcal{M}$。然后
证明 $(\mathcal{M},\varphi)$ 延拓为 $\widehat{S}'$ 上的二元组
$(\mathcal{M}',\varphi')$。当然，此证明与基于 Jouanolou 技巧的证明
本质等价。
%% end units/ch02_u05.tex
